\chapter{Łańcuchy Markowa w czasie dyskretnym}
\label{ch:markov}
\paragraph{Odsyłacze międzytomowe.}
Uogólnienie łańcuchów na czas ciągły znajduje się w tomie II
\IfFileExists{probability_volume_ii.aux}{(rozdział~\ref{II-ch:ctmc})}{(rozdział o łańcuchach Markowa w czasie ciągłym)},
zaś algorytmy MCMC w tomie IV
\IfFileExists{probability_volume_iv.aux}{(rozdział~\ref{IV-ch:mcmc})}{(rozdział o metodach MCMC)}.


\section{Dlaczego przeszłość można zastąpić stanem?}

Wyobraźmy sobie układ, który każdego dnia znajduje się w jednym
z ustalonych stanów. Aby przewidzieć następny dzień, zwykle
potrzebujemy pewnej informacji o przeszłości. W modelu Markowa
zakładamy, że \emph{aktualny stan zawiera całą informację o
przeszłości potrzebną do przewidywania następnego stanu}.
Nie jest to założenie o niezależności kolejnych stanów: przeciwnie,
stan jutrzejszy zazwyczaj zależy od dzisiejszego. Kluczowe jest
to, że zależność od wcześniejszych stanów znika \emph{po
uwarunkowaniu na stan obecny}.

\begin{definicja}[Jednorodny łańcuch Markowa]
Niech $(\Omega,\mathcal F,\Prob)$ będzie przestrzenią
probabilistyczną, $S$ niepustym zbiorem skończonym lub
przeliczalnym, wyposażonym w $2^S$, oraz
$X_n:\Omega\to S$ ($n\in\N_0$) ciągiem zmiennych losowych.
Mówimy, że $(X_n)_{n\geq0}$ jest \emph{jednorodnym
łańcuchem Markowa} o macierzy przejścia
$P=(p_{ij})_{i,j\in S}$, gdy
\[
 p_{ij}\geq0,\qquad \sum_{j\in S}p_{ij}=1
 \quad(i\in S),
\]
oraz dla każdego $n\geq0$ i $j\in S$ zachodzi
\[
\Prob(X_{n+1}=j\mid\sigma(X_0,\ldots,X_n))
=p_{X_nj}\quad\text{p.n.}
\]
Zapis $\Prob(B\mid\mathcal G)$ oznacza tutaj
$\E[\chi_B\mid\mathcal G]$ (rozdział~\ref{ch:condexp}).
Początkowy rozkład $\mu_i=\Prob(X_0=i)$
razem z $P$ określa wszystkie skończenie wymiarowe
rozkłady procesu, jak wykażemy poniżej. Macierz $P$ nazywamy
\emph{stochastyczną wierszowo}.
\end{definicja}

\begin{uwaga}
Gdy $\Prob(X_n=i)=0$, zapis
$\Prob(X_{n+1}=j\mid X_n=i)$ jako iloraz liczb
nie jest zdefiniowany. Sformułowanie przez
warunkową wartość oczekiwaną nie ma tej wady:
zmienna $p_{X_nj}$ jest określona na całej $\Omega$,
a równość z warunkową wartością oczekiwaną zachodzi
prawie na pewno. Wiersz $i$ macierzy $P$ nie musi być
jednoznacznie wyznaczony przez rozkład procesu, jeśli
$\Prob(X_n=i)=0$ dla każdego $n$. Jeśli przejścia zależą
od $n$, otrzymujemy łańcuch niejednorodny
z macierzami $P_n$.
\end{uwaga}


\begin{uwaga}[Jak czytać prawdopodobieństwa dotyczące procesu?]
Zmienna $X_k$ opisuje stan w chwili $k$; zapis
$X_k=i_k$ jest \emph{zdarzeniem} $\{\omega:X_k(\omega)=i_k\}$.
Zapis
\[
\Prob(X_0=i_0,\ldots,X_n=i_n)
=\Prob\left(\bigcap_{k=0}^n\{X_k=i_k\}\right)
\]
oznacza prawdopodobieństwo \emph{dokładnie wskazanej historii},
a nie dowolnej historii, która kończy się stanem $i_n$.
Przecinek wewnątrz argumentu $\Prob$ oznacza więc
przecięcie zdarzeń. Natomiast
$\Prob(X_n=j)$ sumuje prawdopodobieństwa \emph{wszystkich}
historii kończących się w $j$.

Podobnie zapis
$\Prob(X_{n+1}=j\mid X_0=i_0,\ldots,X_n=i_n)$
oznacza iloraz prawdopodobieństw zdarzeń, pod warunkiem
że wskazana historia ma dodatnie prawdopodobieństwo.
Pisanie $\Prob(X_0\mid X_1=j)$ bez wskazania zdarzenia
związanego z $X_0$ nie oznacza jednej liczby.
Można natomiast mówić o \emph{rozkładzie warunkowym}
$\Law(X_0\mid X_1=j)$, podając prawdopodobieństwa
$\Prob(X_0=i\mid X_1=j)$ dla wszystkich $i\in S$,
lub o liczbie $\E[X_0\mid X_1=j]$, gdy stany są liczbami
i $X_0$ jest całkowalna. Są to różne obiekty.
\end{uwaga}

\begin{stwierdzenie}[Wersja elementarna własności Markowa]
Niech $(X_n)$ będzie jednorodnym łańcuchem Markowa
na przeliczalnym $S$ o macierzy $P$. Jeżeli
$\Prob(X_0=i_0,\ldots,X_n=i_n)>0$, to dla każdego $j\in S$
\[
\Prob(X_{n+1}=j\mid X_0=i_0,\ldots,X_n=i_n)=p_{i_nj}.
\]
Niezależnie od pierwszej hipotezy, jeśli
$\Prob(X_n=i)>0$, to
$\Prob(X_{n+1}=j\mid X_n=i)=p_{ij}$.
\end{stwierdzenie}
\begin{proof}
Niech $A=\{X_0=i_0,\ldots,X_n=i_n\}\in
\sigma(X_0,\ldots,X_n)$. Z definicji łańcucha
$\E[\chi_{\{X_{n+1}=j\}}\mid
\sigma(X_0,\ldots,X_n)]=p_{X_nj}$.
Mnożymy tę równość przez $\chi_{A}$ i bierzemy
wartości oczekiwane. Otrzymujemy
$\Prob(A\cap\{X_{n+1}=j\})=p_{i_nj}\Prob(A)$.
Dzielenie przez dodatnie $\Prob(A)$ daje pierwszą
równość. Druga wynika analogicznie, wybierając
$A=\{X_n=i\}$. Jeżeli $\Prob(A)=0$, dzielenie
nie jest dozwolone: wtedy pozostaje poprawne
sformułowanie przez warunkową wartość oczekiwaną.
\end{proof}

\begin{uwaga}[Co jest założeniem, a co twierdzeniem?]
Własności Markowa nie można udowodnić dla \emph{dowolnego}
procesu losowego: definiuje ona szczególną klasę procesów.
Można natomiast udowodnić, że konkretny proces ma tę własność,
a następnie wyprowadzić z niej prawo trajektorii,
własności wielokrokowe i mocną własność Markowa.
Osobne znaczenie ma \emph{zanikanie wpływu stanu początkowego}
po wielu krokach. Jest to twierdzenie asymptotyczne,
którego wystarczającymi warunkami dla skończonych
łańcuchów są nieprzywiedlność i nieokresowość
(nie są one konieczne dla każdego możliwego modelu).
\end{uwaga}

\begin{twierdzenie}[Prawo trajektorii]
Jeżeli $(X_n)$ jest jednorodnym łańcuchem Markowa
na $S$ o rozkładzie początkowym $\mu$ i macierzy $P$,
to dla każdego $n\geq0$ i $i_0,\ldots,i_n\in S$
\[
\Prob(X_0=i_0,\ldots,X_n=i_n)
=\mu_{i_0}\prod_{k=1}^n p_{i_{k-1}i_k},
\]
przy czym iloczyn pusty dla $n=0$ jest równy $1$.
\end{twierdzenie}
\begin{proof}
Dla $n=0$ równość jest definicją $\mu$.
Dla $n\geq1$ załóżmy wzór dla $n-1$ i niech
$A=\{X_0=i_0,\ldots,X_{n-1}=i_{n-1}\}$.
Ponieważ $\chi_{A}$ jest $\sigma(X_0,\ldots,X_{n-1})$-
mierzalna, z definicji łańcucha i własności
wyciągania czynnika otrzymujemy
\[
\Prob(A\cap\{X_n=i_n\})
=\E\left[\chi_{A}\E[\chi_{\{X_n=i_n\}}
\mid\sigma(X_0,\ldots,X_{n-1})]\right]
=p_{i_{n-1}i_n}\Prob(A).
\]
Indukcja daje wymagany iloczyn, również wtedy,
gdy $\Prob(A)=0$; nie dzielimy tu przez
prawdopodobieństwo zdarzenia warunkującego.
\end{proof}


\begin{twierdzenie}[Iloczynowe prawo trajektorii charakteryzuje łańcuch]
Niech $S$ będzie przeliczalny, $\mu$ rozkładem na $S$,
$P$ macierzą stochastyczną, a $(X_n)_{n\geq0}$ procesem
o wartościach w $S$. Załóżmy, że dla wszystkich $n\geq0$
i wszystkich $i_0,\ldots,i_n\in S$ zachodzi
\[
\Prob(X_0=i_0,\ldots,X_n=i_n)
=\mu_{i_0}\prod_{k=1}^np_{i_{k-1}i_k}.
\]
Wtedy proces jest jednorodnym łańcuchem Markowa
o rozkładzie początkowym $\mu$ i macierzy $P$.
\end{twierdzenie}
\begin{proof}
Dla historii $A=\{X_0=i_0,\ldots,X_n=i_n\}$
o dodatnim prawdopodobieństwie dzielenie wzoru
dla $n+1$ przez wzór dla $n$ daje
$\Prob(X_{n+1}=j\mid A)=p_{i_nj}$.
Każde zdarzenie z $\sigma(X_0,\ldots,X_n)$
jest sumą rozłącznych atomów postaci $A$,
ponieważ $S^{n+1}$ jest przeliczalny.
Na atomach o prawdopodobieństwie zero obie całki
są zerowe. Na pozostałych otrzymana równość mówi,
że $p_{X_nj}$ spełnia definicyjną tożsamość całkową.
Funkcja ta jest mierzalna względem
$\sigma(X_0,\ldots,X_n)$ i ograniczona przez $1$.
Sumowanie tożsamości po rozłącznych atomach
jest więc uprawnione i daje ją dla każdego zdarzenia
z tego $\sigma$-ciała.
Wobec tego
$\E[\chi_{\{X_{n+1}=j\}}\mid
\sigma(X_0,\ldots,X_n)]=p_{X_nj}$ p.n.
To jest własność Markowa.
\end{proof}

\begin{twierdzenie}[Warunkowanie wstecz: przeszłość i przyszłość]
Niech $(X_n)$ będzie jednorodnym łańcuchem Markowa
na przeliczalnym $S$ o rozkładzie początkowym
$\mu$ i macierzy $P$ oraz niech $n\geq1$.
Jeśli historia $\{X_1=i_1,\ldots,X_n=i_n\}$
ma dodatnie prawdopodobieństwo, to dla każdego $i\in S$
\[
\Prob(X_0=i\mid X_1=i_1,\ldots,X_n=i_n)
=\Prob(X_0=i\mid X_1=i_1)
=\frac{\mu_i p_{ii_1}}{\sum_{k\in S}\mu_kp_{ki_1}}.
\]
W szczególności, po poznaniu $X_1$ dalsze stany
$X_2,\ldots,X_n$ nie dostarczają dodatkowej informacji
o $X_0$. Nie należy mylić tego z niezależnością
$X_0$ i $X_1$.
\end{twierdzenie}
\begin{proof}
Ze wzoru na prawo trajektorii
\begin{align*}
\Prob(X_0=i, X_1=i_1,\ldots,X_n=i_n)
&=\mu_i p_{ii_1}\,c,\\
\Prob(X_1=i_1,\ldots,X_n=i_n)
&=\left(\sum_{k\in S}\mu_kp_{ki_1}\right)c,
\end{align*}
gdzie $c=\prod_{r=2}^np_{i_{r-1}i_r}$,
a dla $n=1$ przyjmujemy $c=1$. Prawdopodobieństwo
warunkujące jest dodatnie, więc $c>0$ i
$\sum_k\mu_kp_{ki_1}>0$. Po podzieleniu
obu wzorów wspólny czynnik $c$ się skraca.
Tak samo przy $n=1$ otrzymujemy drugi warunek.
Aby przejść do warunkowej niezależności, ustalmy
$j\in S$ z $\Prob(X_1=j)>0$ i dowolną historię
$B=\{X_2=i_2,\ldots,X_n=i_n\}$. Jeśli
$\Prob(B\mid X_1=j)>0$, wykazana równość daje
\[
\Prob(X_0=i,B\mid X_1=j)
=\Prob(X_0=i\mid X_1=j)\Prob(B\mid X_1=j).
\]
Przy zerowym prawdopodobieństwie warunkowym $B$
obie strony są zerowe. Sumowanie po przeliczalnych
zbiorach stanów i historii rozszerza tożsamość
na dowolne zdarzenia z $\sigma(X_0)$
i $\sigma(X_2,\ldots,X_n)$. To jest warunkowa
niezależność względem $\sigma(X_1)$; na jego atomach
zerowych równość jest wymagana tylko p.n.
\end{proof}


\begin{uwaga}[Dlaczego można warunkować zarówno do przodu, jak i wstecz?]
Dla trzech chwil nie należy pisać
$\Prob(X_0\mid X_1,X_2)$ jako liczby.
Jeżeli $x_1,x_2$ są zaobserwowane, szukamy np.
$\Prob(X_0=i\mid X_1=x_1,X_2=x_2)$.
Po ustaleniu $X_1=x_1$ przyszły czynnik
$p_{x_1x_2}$ jest \emph{ten sam} dla każdego
kandydata $X_0=i$, więc skraca się w ilorazie.
Jest to szczególny przypadek warunkowej niezależności:
$\sigma(X_0)$ oraz $\sigma(X_2,\ldots,X_n)$
są warunkowo niezależne względem $\sigma(X_1)$.
Bez poznania $X_1$ taka niezależność nie musi zachodzić.
\end{uwaga}

\begin{uwaga}[Łańcuch odwrócony w czasie]
Jeśli $\mu_n(j)=\Prob(X_n=j)$, to dla $n\geq1$
i $\mu_n(j)>0$ zachodzi
\[
\Prob(X_{n-1}=i\mid X_n=j)
=\frac{\mu_{n-1}(i)p_{ij}}{\mu_n(j)}.
\]
Wynik jest formułą Bayesa zastosowaną do dwóch
sąsiednich stanów; w ogólności zależy od $n$,
ponieważ rozkłady $\mu_{n-1}$ i $\mu_n$ mogą się zmieniać.
Jeśli $\mu_0=\pi$ jest stacjonarny, macierz
przejścia wstecz ma stałe współrzędne
$\widetilde p_{ji}=\pi_i p_{ij}/\pi_j$
dla $\pi_j>0$. Jeśli $\pi_j=0$, wiersz wsteczny
można wybrać jako dowolny rozkład na $S$, ponieważ
taki stan ma prawdopodobieństwo zero w każdej chwili.
Równowaga szczegółowa oznacza
$\widetilde p_{ji}=p_{ji}$ dla $\pi_j>0$.
Wynika to z równości $\pi_i p_{ij}=\pi_j p_{ji}$.
\end{uwaga}

\begin{przyklad}[Pogoda z pamięcią jednego dnia]
Niech $S=\{s,d\}$ oznacza odpowiednio dzień suchy
i deszczowy, a
\[
P=\begin{pmatrix} 4/5&1/5\\ 2/5&3/5\end{pmatrix}.
\]
Jeżeli dzisiaj jest sucho, szansa deszczu jutro
wynosi $1/5$. Prawdopodobieństwo sekwencji
$s,d,d,s$, jeśli zaczynamy na pewno od $s$, jest
$(1/5)(3/5)(2/5)=6/125$. Nie mnożymy
bezwarunkowych prawdopodobieństw dni: kolejne stany
nie są niezależne.
\end{przyklad}


\begin{figure}[htbp]
\centering
\begin{tikzpicture}[>=Stealth,
  stan/.style={draw,very thick,ellipse,minimum width=2.5cm,
  minimum height=1.2cm,align=center,font=\small}]
\node[stan,draw=blue!65!black,fill=blue!8] (s) at (0,0) {sucho\\$s$};
\node[stan,draw=orange!75!black,fill=orange!13] (d) at (6,0) {deszcz\\$d$};
\draw[->,thick,blue!65!black] (s) edge[loop above,looseness=5] node[above] {$4/5$} (s);
\draw[->,thick,orange!75!black] (d) edge[loop above,looseness=5] node[above] {$3/5$} (d);
\draw[->,thick] (s) edge[bend left=16] node[above] {$1/5$} (d);
\draw[->,thick] (d) edge[bend left=16] node[below] {$2/5$} (s);
\end{tikzpicture}
\caption{Graf skierowany dwustanowego łańcucha pogody.
Każda strzałka odpowiada jednej możliwości przejścia;
suma etykiet strzałek wychodzących z danego stanu wynosi $1$.
Pętla oznacza pozostanie w tym samym stanie, a nie brak kolejnego losowania.}
\label{fig:rp1-12-weather-chain}
\end{figure}

\begin{przyklad}[Jedna ścieżka, wiele możliwych ścieżek]
Dla przedstawionego łańcucha i $X_0=s$ p.n. mamy
$\Prob(X_0=s,X_1=d,X_2=d,X_3=s)=6/125$.
Natomiast $\Prob(X_2=d\mid X_0=s)$ jest sumą
prawdopodobieństw dwóch dopuszczalnych historii
$(s,s,d)$ i $(s,d,d)$:
\[
\Prob(X_2=d\mid X_0=s)
=\underbrace{\frac45\frac15}_{s\to s\to d}
+\underbrace{\frac15\frac35}_{s\to d\to d}
=\frac7{25}.
\]
Drzewo rozważanych trajektorii przedstawia
rysunek~\ref{fig:rp1-12-weather-paths}, natomiast
rysunek~\ref{fig:rp1-12-weather-chain} przedstawia
samą macierz przejścia.
To wyjaśnia różnicę między \emph{prawdopodobieństwem
trajektorii} i \emph{prawdopodobieństwem stanu w danej chwili}.
Przykładowo
$\Prob(X_2=d\mid X_0=s,X_1=d)=3/5$,
ponieważ po poznaniu stanu $X_1$ historia wcześniejsza
nie zmienia prognozy następnego kroku.
\end{przyklad}


\begin{figure}[htbp]
\centering
\begin{tikzpicture}[>=Stealth,scale=.94,
  every node/.style={font=\small},
  pathstate/.style={circle,draw=blue!65!black,fill=blue!8,
  minimum size=9mm,inner sep=1pt}]
\node[pathstate] (s0) at (0,0) {$s$};
\node[pathstate] (s1) at (3,1.45) {$s$};
\node[pathstate,draw=orange!80!black,fill=orange!10] (d1) at (3,-1.45) {$d$};
\node[pathstate] (ss) at (6,2.3) {$s$};
\node[pathstate,draw=orange!80!black,fill=orange!10] (sd) at (6,.6) {$d$};
\node[pathstate] (ds) at (6,-.6) {$s$};
\node[pathstate,draw=orange!80!black,fill=orange!10] (dd) at (6,-2.3) {$d$};
\draw[->,thick] (s0)--node[above]{$4/5$}(s1);
\draw[->,thick] (s0)--node[below]{$1/5$}(d1);
\draw[->] (s1)--node[above]{$4/5$}(ss);
\draw[->,very thick,orange!75!black] (s1)--node[below]{$1/5$}(sd);
\draw[->] (d1)--node[above]{$2/5$}(ds);
\draw[->,very thick,orange!75!black] (d1)--node[below]{$3/5$}(dd);
\node at (0,3.05) {$n=0$};\node at (3,3.05) {$n=1$};\node at (6,3.05) {$n=2$};
\end{tikzpicture}
\caption{Drzewo dwóch kroków z początkowego stanu suchego.
W przeciwieństwie do grafu przejść ten rysunek odróżnia
historie. Pomarańczowe strzałki wskazują dwa zakończenia
w stanie deszczowym. Sumujemy prawdopodobieństwa
całych ścieżek do tych zakończeń, a nie same etykiety
ostatnich strzałek. Prawdopodobieństwo pojedynczej
ścieżki jest iloczynem etykiet jej kolejnych krawędzi.}
\label{fig:rp1-12-weather-paths}
\end{figure}

\begin{przyklad}[Proces, który nie ma własności Markowa]
Najpierw losujemy ukrytą zmienną $H\in\{1,2\}$
z prawdopodobieństwami $1/2,1/2$. Jeśli $H=1$,
losowania $X_0,X_1,\ldots$ są niezależnymi próbami
Bernoulliego o parametrze $1/4$; jeśli $H=2$,
mają parametr $3/4$. Bez znajomości $H$
obserwacje nie tworzą łańcucha Markowa.
Rzeczywiście, dla $p_1=1/4$, $p_2=3/4$,
z prawa całkowitego prawdopodobieństwa względem
$H$ i warunkowej niezależności obserwacji przy
ustalonym $H$ otrzymujemy (w obu ilorazach
skraca się wspólny czynnik $1/2$):
\begin{align*}
\Prob(X_2=1\mid X_1=1,X_0=1)
&=\frac{p_1^3+p_2^3}{p_1^2+p_2^2}=\frac7{10},\\
\Prob(X_2=1\mid X_1=1,X_0=0)
&=\frac{p_1^2(1-p_1)+p_2^2(1-p_2)}
{p_1(1-p_1)+p_2(1-p_2)}=\frac12.
\end{align*}
Mimo jednakowego stanu obecnego $X_1=1$ prognoza
zależy od $X_0$. Powód: przeszłe obserwacje
pozwalają wnioskować o ukrytym parametrze $H$,
którego sam $X_1$ nie opisuje. Gdyby do stanu
dołączyć $H$, uzyskalibyśmy proces $(X_n,H)$
o własności Markowa: ze stanu $(x,h)$ przechodzi on
do $(1,h)$ z prawdopodobieństwem $p_h$, a do
$(0,h)$ z prawdopodobieństwem $1-p_h$.
Wynika to z warunkowej niezależności kolejnego
losowania od wcześniejszych przy ustalonym $H=h$.
\end{przyklad}

\begin{przyklad}[Ruch po grafie]
Niech $G=(V,E)$ będzie skończonym nieskierowanym
grafem prostym o niepustym zbiorze wierzchołków,
w którym każdy wierzchołek ma dodatni stopień $d(v)$.
Z wierzchołka $v$ przechodzimy jednostajnie do jednego
z jego sąsiadów: $p_{vw}=1/d(v)$, gdy $\{v,w\}\in E$,
i zero w przeciwnym razie. Jest to spacer losowy
po grafie. Mimo że wybór sąsiada jest jednostajny,
w skończonym spójnym grafie częstości odwiedzin
będą proporcjonalne do stopni wierzchołków,
co wynika z twierdzenia ergodycznego podanego dalej.
Nie zakładamy tutaj zbieżności rozkładów w pojedynczych
chwilach (graf może być dwudzielny).
\end{przyklad}

\section{Macierze przejścia, półgrupa i własność Markowa}

\begin{twierdzenie}[Chapmana--Kołmogorowa]
Niech $(X_n)$ będzie jednorodnym łańcuchem Markowa
o macierzy $P$ na przeliczalnym $S$. Dla $m,n\geq0$
przyjmijmy $p_{ij}^{(n)}=(P^n)_{ij}$, przy czym
$P^0=I$, gdzie $I$ jest macierzą jednostkową.
Dla macierzy stochastycznych $A,B$ na $S$
iloczyn określamy przez $(AB)_{ij}=\sum_k a_{ik}b_{kj}$.
Szeregi są nieujemne i ograniczone przez $1$, a
$\sum_j(AB)_{ij}=\sum_k a_{ik}\sum_j b_{kj}=1$.
Łączność mnożenia wynika z zamiany kolejności
sumowania nieujemnego szeregu podwójnego.
W szczególności potęgi $P^n$ są stochastyczne.
Wówczas
\[
\Prob(X_{m+n}=j\mid\sigma(X_0,\ldots,X_m))
=p^{(n)}_{X_mj}\quad\text{p.n.},\qquad
p_{ij}^{(m+n)}=\sum_{k\in S}p_{ik}^{(m)}p_{kj}^{(n)}.
\]
\end{twierdzenie}
\begin{proof}
Oznaczmy $\mathcal F_m^X=\sigma(X_0,\ldots,X_m)$.
Pokażemy indukcyjnie względem $n$, że
$\E[\chi_{\{X_{m+n}=j\}}\mid\mathcal F_m^X]
=(P^n)_{X_mj}$. Dla $n=0$ jest to
$\chi_{\{X_m=j\}}=(P^0)_{X_mj}$, a dla $n=1$
równość jest właśnie własnością Markowa.
Załóżmy teraz tezę dla $n\geq1$ i każdego
początkowego czasu. Zastosujmy ją z czasem $m+1$,
a następnie użyjmy własności wieży:
\begin{align*}
&\E[\chi_{\{X_{m+n+1}=j\}}\mid\mathcal F_m^X]\\
&=\E[(P^n)_{X_{m+1}j}\mid\mathcal F_m^X]\\
&=\E\left[\sum_{k\in S}(P^n)_{kj}
 \chi_{\{X_{m+1}=k\}}\mid\mathcal F_m^X\right]\\
&=\sum_{k\in S}(P^n)_{kj}\,p_{X_mk}
 =(P^{n+1})_{X_mj}.
\end{align*}
Przy przedostatniej równości wyrazy szeregu są
nieujemne, a ich sumy częściowe są ograniczone
przez $1$ (w każdym punkcie co najwyżej jeden
wyraz jest niezerowy). Można zatem przejść z sumą
pod znak warunkowej wartości oczekiwanej przez
warunkową zbieżność monotoniczną
(rozdział~\ref{ch:condexp}). Ostatnią równość daje
definicja mnożenia macierzy stochastycznych.

Aby odczytać z tego wynik probabilistyczny,
wystarczy wziąć $A=\{X_m=i\}$ o dodatnim
prawdopodobieństwie i pomnożyć otrzymaną
tożsamość przez $\chi_{A}$; dostajemy
$\Prob(X_{m+n}=j\mid X_m=i)=(P^n)_{ij}$.
Nie korzystamy z takiego ilorazu, gdy $\Prob(A)=0$.

Pozostaje równość Chapmana--Kołmogorowa. Z
identyczności macierzowej $P^{m+n}=P^mP^n$
wynika dla wszystkich $i,j$:
$(P^{m+n})_{ij}=\sum_{k\in S}(P^m)_{ik}(P^n)_{kj}$.
Alternatywnie można to uzyskać, sumując po stanie
pośrednim $X_m=k$ dla łańcucha skonstruowanego
z deterministycznym $X_0=i$. Ten drugi opis
uzasadnia intuicję: każda droga długości $m+n$
przechodzi po $m$ krokach przez dokładnie jeden
stan $k$, więc rozłączne możliwości należy zsumować.
\end{proof}


\begin{twierdzenie}[Pełna własność zapominania przeszłości]
Niech $(X_n)_{n\geq0}$ będzie jednorodnym łańcuchem
Markowa na przeliczalnym zbiorze $S$, o macierzy $P$,
a $\mathcal F_n^X=\sigma(X_0,\ldots,X_n)$. Dla
$n\geq0$, $r\geq1$ i ograniczonej funkcji
$f:S^r\to\R$ mamy
\[
\E[f(X_{n+1},\ldots,X_{n+r})\mid\mathcal F_n^X]
=\sum_{j_1,\ldots,j_r\in S} f(j_1,\ldots,j_r)
 p_{X_nj_1}p_{j_1j_2}\cdots p_{j_{r-1}j_r}
\quad\text{p.n.}
\]
W szczególności, po poznaniu $X_n$ rozkład
\emph{całego skończonego fragmentu przyszłości}
nie zależy od wcześniejszych stanów $X_0,\ldots,X_{n-1}$.
\end{twierdzenie}
\begin{proof}
Wystarczy najpierw rozważyć funkcję charakterystyczną
$f=\chi_{\{(j_1,\ldots,j_r)\}}$. Ustalmy historię
$A=\{X_0=i_0,\ldots,X_n=i_n\}$ o dodatnim
prawdopodobieństwie. Z prawa trajektorii
\begin{align*}
&\Prob(A, X_{n+1}=j_1,\ldots,X_{n+r}=j_r)\\
&\qquad=\Prob(A)\,p_{i_nj_1}p_{j_1j_2}\cdots
p_{j_{r-1}j_r}.
\end{align*}
Dzielenie przez $\Prob(A)$ daje odpowiednią
warunkową równość. Jeżeli $\Prob(A)=0$, obie
strony tożsamości całkowej na $A$ mają wartość zero.
Zdarzenia $A$ tworzą przeliczalny podział $\Omega$,
a każde zdarzenie z $\mathcal F_n^X$ jest sumą
niektórych elementów tego podziału. Zsumowanie
równości po tych elementach pokazuje, że prawa
strona jest wersją wskazanej warunkowej wartości
oczekiwanej. Dla dowolnej ograniczonej $f$ mnożymy
równości dla funkcji charakterystycznych przez
$f(j_1,\ldots,j_r)$ i sumujemy po wszystkich ciągach
stanów. Suma iloczynów prawdopodobieństw
trajektorii od ustalonego $i$ wynosi $1$ przez
kolejne sumowania od ostatniego stanu.
Dlatego szereg po prawej stronie jest absolutnie
zbieżny, mierzalny względem $\sigma(X_n)$
i ograniczony modułem przez $\|f\|_\infty$.
Po lewej stronie sumy częściowe funkcji
$f(X_{n+1},\ldots,X_{n+r})$ są również ograniczone
przez $\|f\|_\infty$, gdyż zdarzenia są rozłączne.
Warunkowa zbieżność zmajoryzowana uzasadnia
przejście do granicy. Wyjątki dla przeliczalnie
wielu równości można zebrać w jeden zbiór zerowy.
\end{proof}

\begin{uwaga}[Dwa różne znaczenia słowa „zapominanie”]
Powyższa własność jest \emph{dokładna i natychmiastowa}:
przeszłość przestaje mieć znaczenie dla dalszej prognozy,
gdy stan obecny jest znany. Nie przesądza ona o tym,
czy $\Prob(X_n=j\mid X_0=i)$ zbliża się, gdy
$n\to\infty$, do liczby niezależnej od $i$.
Takie \emph{asymptotyczne zapominanie początku}
wynika w skończonym przypadku z nieprzywiedlności
i nieokresowości. Dla łańcucha naprzemiennego
$0\leftrightarrow1$ własność Markowa zachodzi,
a pamięć o stanie początkowym pozostaje na zawsze.
\end{uwaga}

Dla ograniczonej $f:S\to\R$ oznaczmy
$(Pf)(i)=\sum_jp_{ij}f(j)$ oraz
$(P^nf)(i)=\sum_jp_{ij}^{(n)}f(j)$.
Sumując równości twierdzenia Chapmana--Kołmogorowa
po $j$ z wagami $f(j)$ i korzystając z warunkowej
zbieżności zmajoryzowanej, otrzymujemy
$\E[f(X_{m+n})\mid\sigma(X_0,\ldots,X_m)]
=(P^nf)(X_m)$. Operator $P$ wyznacza zatem
warunkowe wartości oczekiwane funkcji przyszłych stanów,
natomiast rozkład wierszowy $\mu$ propaguje się jako
$\mu P^n$.

\begin{przyklad}[Dwa stany: rozwiązanie jawne]
Dla $0<\alpha,\beta\leq1$ weźmy
$P=\begin{psmallmatrix}1-\alpha&\alpha\\
\beta&1-\beta\end{psmallmatrix}$.
Jeśli $q_n=\Prob(X_n=1)$, to
$q_{n+1}=\alpha(1-q_n)+(1-\beta)q_n
=\alpha+(1-\alpha-\beta)q_n$.
Po odjęciu $q_*=\alpha/(\alpha+\beta)$ otrzymujemy
$q_{n+1}-q_*=(1-\alpha-\beta)(q_n-q_*)$,
więc
\[
q_n=\frac{\alpha}{\alpha+\beta}
+(1-\alpha-\beta)^n
\left(q_0-\frac{\alpha}{\alpha+\beta}\right).
\]
Przy $\alpha=\beta=1$ łańcuch zmienia stan
w każdym kroku. Wtedy $q_n=1/2+(-1)^n(q_0-1/2)$,
więc rozkłady zbiegają wtedy i tylko wtedy,
gdy $q_0=1/2$ (rozkład od początku jest
stacjonarny). W szczególności nie zbiegają
przy starcie z ustalonego stanu. W pozostałych
przypadkach $\alpha,\beta>0$ i
$\alpha+\beta<2$ mamy $|1-\alpha-\beta|<1$,
więc rozkład dąży do stacjonarnego.
\end{przyklad}


\begin{przyklad}[Ilościowe zapominanie początku a własność Markowa]
Dla łańcucha pogody
$P=\begin{psmallmatrix}4/5&1/5\\2/5&3/5\end{psmallmatrix}$
rozkład stacjonarny jest równy
$\pi=(2/3,1/3)$. Wzór dla łańcucha dwustanowego
daje
\begin{align*}
p_{sd}^{(n)}&=\frac13\bigl(1-(2/5)^n\bigr),\\
p_{dd}^{(n)}&=\frac13+\frac23(2/5)^n.
\end{align*}
Różnica obu prognoz wynosi $(2/5)^n$,
a zatem zależność od stanu początkowego zanika
wykładniczo. Jest to \emph{inny fakt} niż to,
że do prognozowania $X_{n+1}$ wystarcza $X_n$.
Pierwszy mówi o coraz słabszym wpływie dawnej historii;
drugi — o warunkowej niezależności w jednym kroku.
\end{przyklad}

\section{Konstrukcja i łączenie z niezależnymi losowaniami}

Zasadnicze jest pytanie, czy każda macierz $P$ i rozkład
początkowy rzeczywiście wyznaczają łańcuch. Można to
zrobić bez polegania na nieformalnych warunkowych
losowaniach.

\begin{twierdzenie}[Konstrukcja z jednostajnych zmiennych losowych]
Dla dowolnego przeliczalnego $S$, macierzy stochastycznej
$P$ na $S$ oraz rozkładu $\mu$ na $S$ istnieje
jednorodny łańcuch Markowa $(X_n)_{n\geq0}$
o rozkładzie początkowym $\mu$ i macierzy $P$.
\end{twierdzenie}
\begin{proof}
Wybierzmy enumerację $S=\{s_1,s_2,\ldots\}$
(w przypadku skończonym kończymy listę).
Dla $i\in S$ i $u\in(0,1)$ niech $F(i,u)=s_r$,
gdzie $r$ jest najmniejszym indeksem takim, że
$\sum_{\ell\leq r}p_{i,s_\ell}\geq u$.
Definicja ma sens dla każdego $u\in(0,1)$, ponieważ
sumy częściowe dążą do $1$. Oznaczając
$c_{i,r}=\sum_{\ell\leq r}p_{i,s_\ell}$ oraz
$c_{i,0}=0$, mamy dokładnie
\[
F^{-1}(\{s_r\})=
\bigcup_{i\in S}\{i\}\times
\bigl((c_{i,r-1},c_{i,r}]\cap(0,1)\bigr).
\]
Jest to przeliczalna suma zbiorów mierzalnych,
więc $F$ jest mierzalna względem
$2^S\otimes\mathcal B((0,1))$. Długość wskazanego
przedziału wynosi $p_{i,s_r}$, również gdy jest zerowa.
Na przestrzeni produktowej
$S\times(0,1)^{\N}$ z miarą produktową
$\mu\otimes\lambda_1^{\otimes\N}$ (każdy czynnik
Lebesgue'a ograniczamy do $(0,1)$) weźmy
współrzędne $U_1,U_2,\ldots\sim\Unif(0,1)$,
niezależne od siebie i od
$X_0$ o rozkładzie $\mu$, i rekurencyjnie połóżmy
$X_{n+1}=F(X_n,U_{n+1})$. Ponieważ
$U_{n+1}$ jest niezależna od $X_0,U_1,\ldots,U_n$,
a każda $X_k$ dla $k\leq n$ jest mierzalna względem
$\sigma(X_0,U_1,\ldots,U_n)$ na mocy rekurencji,
$U_{n+1}$ jest niezależna od $\sigma(X_0,\ldots,X_n)$.
Dla dowolnego $A$ z tego $\sigma$-ciała mamy
\begin{align*}
\Prob(A\cap\{X_{n+1}=j\})
&=\sum_{i\in S}\Prob(A\cap\{X_n=i\}
                      \cap\{F(i,U_{n+1})=j\})\\
&=\sum_{i\in S}\Prob(A\cap\{X_n=i\})p_{ij}\\
&=\E[\chi_A\,p_{X_nj}].
\end{align*}
Prawa strona jest całką ograniczonej mierzalnej
funkcji $p_{X_nj}$ po $A$. Z definicji warunkowej
wartości oczekiwanej otrzymujemy własność Markowa,
także na atomach o prawdopodobieństwie zero.
\end{proof}

Dalej $\Prob_\mu$ oznacza rozkład trajektorii
na $S^{\N_0}$ wyznaczony przez rozkład początkowy
$\mu$ oraz ustaloną macierz $P$, a $\E_\mu$
odpowiadającą wartość oczekiwaną.
Dla $\mu=\delta_i$ piszemy $\Prob_i$ i $\E_i$.
Prawo to istnieje na mocy powyższej konstrukcji
i jest jednoznaczne, gdyż prawo trajektorii
wyznacza miarę na wszystkich cylindrach, które
generują produktowe $\sigma$-ciało
(rozdział~\ref{ch:probability}).
W tej notacji $X_n$ oznacza $n$-tą współrzędną
trajektorii.

\section{Czasy wejścia, mocna własność Markowa}

\begin{definicja}[Czas wejścia i powrotu]
Dla łańcucha $(X_n)$ i zbioru $A\subseteq S$ definiujemy
$\tau_A=\inf\{n\geq0:X_n\in A\}$ oraz
$\tau_A^+=\inf\{n\geq1:X_n\in A\}$,
przyjmując $\inf\varnothing=\infty$.
Dla jednego stanu piszemy $\tau_i$ i $\tau_i^+$.
\end{definicja}

Czasem zatrzymania względem naturalnej filtracji
$\mathcal F_n^X=\sigma(X_0,\ldots,X_n)$ nazywamy funkcję
$\tau:\Omega\to\N_0\cup\{\infty\}$, dla której
$\{\tau\leq n\}\in\mathcal F_n^X$ dla każdego $n$.
Oba powyższe czasy są czasami zatrzymania, ponieważ
$\{\tau_A\leq n\}=\bigcup_{k=0}^n\{X_k\in A\}$,
a dla $\tau_A^+$ bierzemy sumę od $k=1$ (pustą dla $n=0$).
Dalsze własności czasów zatrzymania rozwiniemy
w rozdziale~\ref{ch:filtracje}.

\begin{twierdzenie}[Mocna własność Markowa, wersja dyskretna]
Niech $(X_n)$ będzie jednorodnym łańcuchem Markowa
na przeliczalnym $S$, a $\tau$ czasem zatrzymania
względem jego naturalnej filtracji
$\mathcal F_n^X=\sigma(X_0,\ldots,X_n)$.
Dla każdego $r\geq0$ i ograniczonej funkcji
$f:S^{r+1}\to\R$ zachodzi
\[
\E\bigl[f(X_\tau,\ldots,X_{\tau+r})
\chi_{\{\tau<\infty\}}\mid\mathcal F_\tau^X\bigr]
=\chi_{\{\tau<\infty\}}\,G_f(X_\tau),
\]
gdzie $G_f(i)=\E_i f(X_0,\ldots,X_r)$, a
$\E_i$ oznacza wartość oczekiwaną dla łańcucha
rozpoczynającego w $i$. Oba wyrażenia przed warunkowaniem i po jego prawej
stronie definiujemy kawałkami: na $\{\tau=\infty\}$
przyjmują wartość zero, bez potrzeby określania
$X_\infty$ (zero nie musi należeć do $S$).
\end{twierdzenie}
\begin{proof}
Przyjmujemy
$\mathcal F_\tau^X=\{A\in\mathcal F:
A\cap\{\tau\leq n\}\in\mathcal F_n^X
\text{ dla każdego }n\geq0\}$.
Jest to $\sigma$-ciało (przecięcie ze
$\{\tau\leq n\}$ zachowuje przeliczalne sumy
i względne dopełnienia). Równoważnie
$A\cap\{\tau=n\}\in\mathcal F_n^X$ dla każdego $n$:
w jedną stronę odejmujemy
$A\cap\{\tau\leq n-1\}$, a w drugą sumujemy
po $k\leq n$; przypadek $n=0$ jest bezpośredni.
Szczegóły omawiamy w rozdziale~\ref{ch:filtracje}.
Dla dowolnego takiego $A$ piszemy, korzystając
z rozbicia $\{\tau<\infty\}=\bigcup_{n\geq0}\{\tau=n\}$,
\begin{align*}
&\E\bigl[\chi_{A}\chi_{\{\tau<\infty\}}
 f(X_\tau,\ldots,X_{\tau+r})\bigr]\\
&=\sum_{n\geq0}\E\bigl[
\chi_{A\cap\{\tau=n\}}f(X_n,\ldots,X_{n+r})\bigr]\\
&=\sum_{n\geq0}\E\bigl[
\chi_{A\cap\{\tau=n\}}G_f(X_n)\bigr]
=\E\bigl[\chi_{A}\chi_{\{\tau<\infty\}}G_f(X_\tau)\bigr].
\end{align*}
Druga równość wynika z pełnej własności Markowa
dla skończonych fragmentów przyszłości:
dla każdego ustalonego $i$ stosujemy ją do
$(j_1,\ldots,j_r)\mapsto f(i,j_1,\ldots,j_r)$,
a następnie sumujemy po rozłącznych zdarzeniach
$\{X_n=i\}$. Dla $r=0$ równość jest bezpośrednia.
Daje to
$\E[f(X_n,\ldots,X_{n+r})\mid\mathcal F_n^X]
=G_f(X_n)$ p.n.
Szeregi są absolutnie zbieżne, bo $f$ jest ograniczona
i zdarzenia $\{\tau=n\}$ są rozłączne.
Prawa strona jest $\mathcal F_\tau^X$-mierzalna:
na $\{\tau=n\}$ jest równa $G_f(X_n)$,
które jest $\mathcal F_n^X$-mierzalne.
Dla borelowskiego zbioru $B$ przeciwobraz prawej
strony przecięty z $\{\tau=n\}$ jest więc równy
$\{\tau=n\}\cap\{G_f(X_n)\in B\}$.
Ponadto jest ograniczona przez $\|f\|_\infty$,
więc całkowalna. Wykazana tożsamość całkowa
definiuje wskazaną warunkową wartość oczekiwaną. Nie ma tu żadnego warunkowania
na zdarzenie $\{\tau=n\}$ o dodatniej miarze:
dowód działa również dla zerowych prawdopodobieństw.
\end{proof}

\begin{przyklad}[Po pierwszym deszczu]
W modelu pogody zatrzymajmy obserwację w pierwszym
dniu deszczowym $\tau_d$. Mocna własność Markowa daje dokładnie
\[
\E[\chi_{\{\tau_d<\infty,\ X_{\tau_d+1}=s\}}
\mid\mathcal F_{\tau_d}^X]
=\frac25\chi_{\{\tau_d<\infty\}}.
\]
Zdarzenie po lewej definiujemy tylko na
skończonych czasach wejścia. W szczególności
dla każdego $n$ z $\Prob(\tau_d=n)>0$ mamy
$\Prob(X_{n+1}=s\mid\tau_d=n)=2/5$.
To właśnie treść mocnej własności Markowa:
czas $\tau_d$ jest losowy, ale przyszłość po
zatrzymaniu znów zaczyna się od obecnego stanu.
\end{przyklad}

\section{Klasy komunikujące, powracanie i okresowość}

\begin{definicja}[Komunikowanie, nieprzywiedlność, okres]
Dla $i,j\in S$ piszemy $i\to j$, jeśli istnieje
$n\geq0$ z $p_{ij}^{(n)}>0$. Stany $i$ i $j$
\emph{komunikują się}, jeśli $i\to j$ i $j\to i$.
Łańcuch jest \emph{nieprzywiedlny}, gdy wszystkie
stany komunikują się. Klasa $C\subseteq S$
jest \emph{domknięta}, jeśli $p_{ij}=0$ dla
$i\in C$, $j\notin C$. Okresem stanu $i$ nazywamy
$d(i)=\gcd\{n\geq1:p_{ii}^{(n)}>0\}$;
przez $\gcd$ rozumiemy największy wspólny dzielnik
niepustego zbioru, zaś jeśli zbiór jest pusty,
stan nie ma dodatnich czasów powrotu i okres
pozostawiamy nieokreślony. Stan o okresie $1$
nazywamy nieokresowym, a łańcuch nieokresowym,
jeśli każdy jego stan ma okres $1$.
\end{definicja}

\begin{lemat}[Okres jest wspólny dla klasy komunikującej]
Jeżeli stany $i,j$ komunikują się i istnieje
$n\geq1$ z $p_{ii}^{(n)}>0$, to $d(i)=d(j)$.
\end{lemat}
\begin{proof}
Dla $i=j$ teza jest oczywista. Dla $i\ne j$
wybierzmy $r,s\geq1$ takie, że
$p_{ij}^{(r)}>0$ i $p_{ji}^{(s)}>0$.
Z równania Chapmana--Kołmogorowa mamy
$p_{ii}^{(r+s)}\geq p_{ij}^{(r)}p_{ji}^{(s)}>0$
oraz $p_{jj}^{(r+s)}>0$, więc oba okresy są określone.
W szczególności $d(i)$ dzieli $r+s$.
Dla dowolnego $n\geq1$ z $p_{jj}^{(n)}>0$ mamy
$p_{ii}^{(r+n+s)}\geq p_{ij}^{(r)}p_{jj}^{(n)}
p_{ji}^{(s)}>0$, zatem $d(i)$ dzieli $r+n+s$
i $r+s$, więc dzieli $n$.
Tym samym $d(i)$ dzieli $d(j)$.
Zamieniając $i$ z $j$, otrzymujemy podzielność
odwrotną, a więc równość.
\end{proof}

\begin{definicja}[Stan powracający i przemijający]
Stan $i\in S$ jest \emph{powracający}, jeśli
$\Prob_i(\tau_i^+<\infty)=1$, a
\emph{przemijający}, jeśli
$\Prob_i(\tau_i^+<\infty)<1$.
Stan powracający jest \emph{dodatnio powracający},
jeżeli $\E_i\tau_i^+<\infty$, a \emph{zerowo
powracający}, jeśli $\E_i\tau_i^+=\infty$.
\end{definicja}

\begin{twierdzenie}[Kryterium powracania przez szereg odwiedzin]
Dla każdego stanu $i$ przeliczalnego jednorodnego
łańcucha Markowa zachodzi równoważność
\[
 i\text{ jest powracający}\quad\Longleftrightarrow
\quad\sum_{n=0}^{\infty}p_{ii}^{(n)}=\infty.
\]
Jeżeli $r_i=\Prob_i(\tau_i^+<\infty)<1$,
to $\sum_{n\geq0}p_{ii}^{(n)}=1/(1-r_i)$.
\end{twierdzenie}
\begin{proof}
Rozważmy liczbę wszystkich wizyt w $i$,
$N_i=\sum_{n\geq0}\chi_{\{X_n=i\}}$, dla łańcucha
rozpoczętego w $i$. Niech $T_0=0$, a rekurencyjnie
$T_{k+1}=\inf\{n>T_k:X_n=i\}$, z wartością $\infty$,
jeśli taki czas nie istnieje. Każdy $T_k$ jest czasem
zatrzymania: dla $k\geq1$ zdarzenie $\{T_k\leq n\}$
jest równe $\{\sum_{\ell=1}^n\chi_{\{X_\ell=i\}}\geq k\}$.
Dla $m\geq1$ mocna własność Markowa, zastosowana
do zdarzenia powrotu w ciągu następnych $m$ kroków,
daje
\[
\Prob_i(T_{k+1}<\infty,\ T_{k+1}-T_k\leq m)
=\Prob_i(T_k<\infty)\Prob_i(\tau_i^+\leq m),
\]
gdzie zdarzenie po lewej rozumiemy tylko dla
skończonego $T_{k+1}$. Przechodząc do granicy
$m\to\infty$ przez ciągłość miary z dołu,
otrzymujemy
$\Prob_i(T_{k+1}<\infty)=r_i\Prob_i(T_k<\infty)$.
Indukcja daje $\Prob_i(T_k<\infty)=r_i^k$.
Ponieważ punktowo $N_i=\sum_{k\geq0}\chi_{\{T_k<\infty\}}$,
dwukrotne zastosowanie twierdzenia o zbieżności monotonicznej daje
\[
\sum_{n\geq0}p_{ii}^{(n)}=\E_iN_i
=\sum_{k\geq0}\Prob_i(T_k<\infty)
=\sum_{k\geq0}r_i^k.
\]
Ostatni szereg jest skończony dokładnie dla $r_i<1$;
przy $r_i=1$ jest rozbieżny. Dowód obejmuje oba
przypadki bez zakładania, że czas powrotu ma
skończoną wartość oczekiwaną.
\end{proof}

\begin{przyklad}[Spacer symetryczny na prostej całkowitej]
Niech $X_0=0$ i $X_n=\xi_1+\cdots+\xi_n$,
gdzie $\xi_k$ są niezależne i przyjmują $\pm1$
z prawdopodobieństwem $1/2$. Dla $n\geq0$
$\Prob_0(X_{2n+1}=0)=0$ i
$\Prob_0(X_{2n}=0)=\binom{2n}{n}2^{-2n}$.
Ze wzoru Stirlinga
$\binom{2n}{n}2^{-2n}\sim1/\sqrt{\pi n}$.
Szereg prawdopodobieństw powrotu jest rozbieżny,
więc spacer jest powracający. Ponadto okres wynosi $2$.
Dla spaceru z niezależnymi przyrostami o wspólnym
rozkładzie $\Prob(\xi_k=1)=p$, gdzie $0<p<1$
i $p\ne1/2$, oraz $\Prob(\xi_k=-1)=1-p$,
mamy
$\Prob_0(X_{2n}=0)=\binom{2n}{n}[p(1-p)]^n
\sim[4p(1-p)]^n/\sqrt{\pi n}$,
a $4p(1-p)<1$; przez porównanie z szeregiem
geometrycznym stan zero jest zatem przemijający.
Dla $p=0$ lub $p=1$ spacer porusza się
deterministycznie w jednym kierunku i również
nie powraca do punktu wyjścia.
Przez niezmienniczość przejść względem przesunięć
te same wnioski o powracaniu dotyczą każdego stanu.
Dla $0<p<1$ (w tym w przypadku symetrycznym)
okres każdego stanu wynosi $2$: powroty
w nieparzystych chwilach są niemożliwe, a powrót
po dwóch krokach ma prawdopodobieństwo $2p(1-p)>0$.
Dla $p=0$ lub $p=1$ okres jest nieokreślony
w przyjętej konwencji.
\end{przyklad}

\begin{przyklad}[Dlaczego przy powracaniu liczy się cały szereg]
Prawdopodobieństwo powrotu w jednym kroku może
być małe, a powrót w pewnej dodatniej chwili może
mimo to zachodzić prawie na pewno,
jak pokazuje kryterium szeregowe.
Sama liczba możliwych chwil powrotu nie wystarcza:
w spacerze z dryfem dodatnie prawdopodobieństwa
powrotu występują w nieskończenie wielu chwilach,
ale ich szereg jest zbieżny. W szczególności
samo $p_{ii}$ nie rozstrzyga o powracaniu.
Na przykład spacer symetryczny ma $p_{00}=0$,
a mimo to stan zero jest powracający.
\end{przyklad}


\begin{stwierdzenie}[Domknięcie i komunikowanie]
Dla jednorodnego łańcucha na przeliczalnym zbiorze $S$
relacja komunikowania jest relacją równoważności.
Jeśli klasa komunikująca $C$ jest domknięta i
$X_0\in C$ p.n., to $X_n\in C$ dla wszystkich
$n\geq0$ prawie na pewno.
\end{stwierdzenie}
\begin{proof}
Zwrotność wynika z $p_{ii}^{(0)}=1$,
a symetryczność jest częścią definicji.
Jeśli $i\to j$ w $m$ krokach i $j\to k$ w $n$,
to z równania Chapmana--Kołmogorowa
$p_{ik}^{(m+n)}\geq p_{ij}^{(m)}p_{jk}^{(n)}>0$.
Tak samo, łącząc drogę powrotną $k\to j\to i$,
dostajemy $k\to i$. To dowodzi przechodniości.
Jeśli $C$ jest domknięta, to dla każdego
$i\in C$ mamy $\Prob_i(X_1\in C)=1$.
Indukcja z warunkowaniem względem $X_n$
daje $\Prob(X_n\in C)=1$ przy $X_0\in C$ p.n.
Przeliczalne przecięcie zdarzeń pełnej miary
zapewnia pozostawanie w $C$ jednocześnie
we wszystkich chwilach.
\end{proof}

\begin{stwierdzenie}[Powracanie jest własnością klasy komunikującej]
Jeżeli $i$ i $j$ komunikują się, to albo oba stany
są powracające, albo oba są przemijające.
Każda klasa powracająca jest domknięta.
Domkniętość klasy nieskończonej nie gwarantuje
jednak powracania, czego przykładem jest spacer
z dryfem na $\Z$.
\end{stwierdzenie}
\begin{proof}
Wybierzmy $r,s\geq0$ z $p_{ji}^{(r)}>0$
i $p_{ij}^{(s)}>0$. Dla każdego $n\geq0$
\[
p_{jj}^{(r+n+s)}\geq p_{ji}^{(r)}
p_{ii}^{(n)}p_{ij}^{(s)}.
\]
Jeżeli szereg po $n$ dla $i$ jest rozbieżny,
to szereg dla $j$ także jest rozbieżny,
a kryterium powracania daje tezę w jedną stronę.
Zamiana $i$ i $j$ daje drugą.
Jeśli klasa powracająca $C$ nie byłaby domknięta,
istniałyby $i\in C$, $j\notin C$ takie, że $p_{ij}>0$.
Wtedy $j$ nie może osiągnąć $i$ (w przeciwnym
razie komunikowałyby się). Zatem
$\Prob_j(\tau_i<\infty)=0$, ponieważ każde
$\Prob_j(X_n=i)=0$. Z własności Markowa
$\Prob_i(X_1=j,\ \tau_i^+=\infty)=p_{ij}>0$,
co przeczy powracaniu $i$.
\end{proof}

\begin{przyklad}[Dwie klasy pochłaniające i stan przejściowy]
Niech $S=\{0,1,2\}$ i
$P=\begin{psmallmatrix}1&0&0\\ 1/3&0&2/3\\ 0&0&1\end{psmallmatrix}$.
Stany $0$ i $2$ są oddzielnymi domkniętymi klasami,
a stan $1$ jest przemijający. Rozkłady
$\delta_0$ i $\delta_2$ są stacjonarne;
stacjonarny jest też każdy ich rozkład mieszany
$t\delta_0+(1-t)\delta_2$, gdzie $t\in[0,1]$.
Są to wszystkie rozkłady stacjonarne: kolumna
macierzy odpowiadająca stanowi $1$ jest zerowa,
więc równanie $\pi P=\pi$ wymusza $\pi_1=0$.
Zatem bez nieprzywiedlności nie musi być
jednoznacznego rozkładu stacjonarnego ani
wspólnej granicy niezależnej od startu.
\end{przyklad}

\begin{uwaga}[Kiedy dowód prawa Markowa przechodzi na proces obserwowany?]
Własność Markowa zależy od wyboru \emph{stanu}.
Jeśli stan rzeczywisty ma dwie składowe $(X_n,H)$,
ale obserwujemy tylko $X_n$, to proces obserwacji
nie musi mieć własności Markowa. Rozkład warunkowy
przyszłości może zależeć od wcześniejszych obserwacji,
gdyż pomagają one odtworzyć ukrytą składową $H$.
W dowodach własności Markowa należy więc wskazać
niezależne innowacje lub zweryfikować tożsamość
warunkową, a nie wnioskować z samej znajomości
jednokrokowych rozkładów bezwarunkowych.
\end{uwaga}

\begin{przyklad}[Rozwiązane zadanie: warunkowanie wstecz]
W modelu pogody z macierzą
$P=\begin{psmallmatrix}4/5&1/5\\2/5&3/5\end{psmallmatrix}$
przyjmijmy $\Prob(X_0=s)=3/4$ i
$\Prob(X_0=d)=1/4$. Wyznaczymy prawdopodobieństwo
obserwacji $X_0=s,X_1=d,X_2=d$ oraz szansę,
że pierwszy dzień był suchy, jeśli widzimy
deszcz w dniach $1$ i $2$.

Z prawa trajektorii
$\Prob(X_0=s,X_1=d,X_2=d)
=(3/4)(1/5)(3/5)=9/100$.
Prawdopodobieństwo warunkujące wynosi
\[
\Prob(X_1=d,X_2=d)
=\left(\frac34\frac15+\frac14\frac35\right)\frac35
=\frac9{50}.
\]
Stąd
$\Prob(X_0=s\mid X_1=d,X_2=d)
=(9/100)/(9/50)=1/2$.
Gdybyśmy znali jedynie $X_1=d$, wynik także
wynosiłby
$[(3/4)(1/5)]/[(3/4)(1/5)+(1/4)(3/5)]=1/2$.
To konkretne sprawdzenie warunkowej niezależności
przeszłości i dalszej przyszłości po poznaniu
stanu pośredniego.
\end{przyklad}

\section{Rozkłady stacjonarne i równowaga}

\begin{definicja}[Rozkład stacjonarny i odwracalność]
Rozkład probabilistyczny $\pi=(\pi_i)_{i\in S}$
jest \emph{stacjonarny} dla $P$, jeśli
$\pi_j=\sum_{i\in S}\pi_i p_{ij}$ dla każdego $j$,
czyli $\pi P=\pi$.
Mówimy, że $\pi$ spełnia \emph{warunek równowagi
szczegółowej}, jeśli
$\pi_i p_{ij}=\pi_jp_{ji}$ dla wszystkich $i,j\in S$.
Łańcuch z rozkładem początkowym $\pi$ spełniającym
równowagę szczegółową nazywamy odwracalnym
względem $\pi$.
\end{definicja}

\begin{stwierdzenie}[Równowaga szczegółowa implikuje stacjonarność]
Jeżeli $\pi$ jest rozkładem probabilistycznym
spełniającym warunek równowagi szczegółowej,
to jest stacjonarny.
\end{stwierdzenie}
\begin{proof}
Dla każdego $j$ sumujemy po $i$:
$\sum_i\pi_i p_{ij}=\sum_i\pi_jp_{ji}
=\pi_j\sum_i p_{ji}=\pi_j$.
W ostatnim wyrażeniu sumujemy \emph{wiersz}
macierzy $P$ o indeksie początkowym $j$.
\end{proof}

Jeżeli $X_0$ ma rozkład stacjonarny $\pi$, to
$X_n$ ma rozkład $\pi P^n=\pi$ dla każdego $n$.
Ponadto z prawa trajektorii i wielokrokowej
własności Markowa, dla $0\leq t_0<\cdots<t_r$
i $i_0,\ldots,i_r\in S$,
\[
\Prob(X_{t_0}=i_0,\ldots,X_{t_r}=i_r)
=\pi_{i_0}\prod_{\ell=1}^r
p_{i_{\ell-1}i_\ell}^{(t_\ell-t_{\ell-1})}.
\]
Wzór nie zmienia się przy wspólnym przesunięciu
chwil o dowolną nieujemną liczbę całkowitą,
co dowodzi ścisłej stacjonarności procesu.
Przy równowadze szczegółowej mamy także
\[
\pi_{i_0}p_{i_0i_1}\cdots p_{i_{r-1}i_r}
=\pi_{i_r}p_{i_ri_{r-1}}\cdots p_{i_1i_0}.
\]
Równość wynika z kolejnego stosowania równowagi
szczegółowej od lewej do prawej, bez dzielenia
przez $\pi_i$; działa zatem również dla zerowych mas.
Wyraża ona niezmienniczość prawa skończonych
trajektorii przy odwróceniu czasu.

\begin{przyklad}[Spacer na grafie: przewaga wierzchołków o dużym stopniu]
Dla skończonego nieskierowanego grafu prostego
o niepustym zbiorze wierzchołków, bez wierzchołków
izolowanych, zdefiniujmy
$\pi_v=d(v)/\sum_{w\in V}d(w)$.
Gdy $v,w$ są sąsiadami,
$\pi_vp_{vw}=1/\sum_zd(z)=\pi_wp_{wv}$;
gdy nie są, obie strony znikają. Zatem $\pi$
jest stacjonarny. Wierzchołek z trzykrotnie
większą liczbą sąsiadów ma w równowadze
trzykrotnie większe prawdopodobieństwo wizyty.
\end{przyklad}

\begin{przyklad}[Stacjonarność nie jest zbieżnością]
Dla $P=\begin{psmallmatrix}0&1\\1&0\end{psmallmatrix}$
rozkład $(1/2,1/2)$ jest stacjonarny, lecz
przy starcie z pierwszego stanu rozkłady
kolejnych $X_n$ naprzemiennie wynoszą
$(1,0)$ i $(0,1)$.
Sama równość $\pi P=\pi$ nie wystarcza do
$\mu P^n\to\pi$; trzeba kontrolować okresowość.
\end{przyklad}

\begin{twierdzenie}[Dodatnie powracanie a rozkład stacjonarny]
\label{thm:rp1-12-positive-recurrence}
Niech $P$ będzie macierzą stochastyczną nieprzywiedlnego
łańcucha Markowa na niepustym przeliczalnym zbiorze $S$.
Następujące warunki są równoważne:
\begin{enumerate}
\item Pewien stan jest dodatnio powracający.
\item Każdy stan jest dodatnio powracający.
\item Istnieje rozkład stacjonarny $\pi$.
\end{enumerate}
Jeżeli te warunki zachodzą, rozkład stacjonarny
jest jedyny i dla każdego $i\in S$ spełnia
\[
\pi_i=\frac1{\E_i\tau_i^+}.
\]
\end{twierdzenie}
\begin{proof}
Ustalmy stan $i$ i oznaczmy $\tau=\tau_i^+$.
Najpierw załóżmy, że $i$ jest dodatnio powracający.
Dla $j\in S$ określmy oczekiwaną liczbę wizyt
w $j$ w jednym cyklu powrotu do $i$:
\[
v_j=\E_i\sum_{n=0}^{\tau-1}\chi_{\{X_n=j\}}.
\]
Mamy $v_i=1$ oraz z twierdzenia Tonellego
$\sum_jv_j=\E_i\tau<\infty$.
Pokażemy, że $vP=v$. Dla każdego $j\in S$
\begin{align*}
\sum_{k\in S}v_kp_{kj}
&=\sum_{n\geq0}\E_i[
      \chi_{\{\tau>n\}}p_{X_nj}]\\
&=\sum_{n\geq0}\Prob_i(\tau>n,X_{n+1}=j)\\
&=\E_i\sum_{n=1}^{\tau}\chi_{\{X_n=j\}}=v_j.
\end{align*}
Druga równość wynika z własności Markowa, ponieważ
$\{\tau>n\}\in\mathcal F_n^X$. Ostatnia zachodzi,
gdyż $X_0=X_\tau=i$ p.n., a zamiany nieujemnych
sum i wartości oczekiwanych uzasadnia twierdzenie
Tonellego. Normalizacja przez $\E_i\tau$ daje
rozkład stacjonarny $\pi_j=v_j/\E_i\tau$.
Otrzymaliśmy implikację (1)$\Rightarrow$(3).

Niech teraz $\pi$ będzie dowolnym rozkładem stacjonarnym.
Z nieprzywiedlności wynika $\pi_i>0$ dla każdego $i$:
wybieramy $k$ z $\pi_k>0$ oraz $m$ z $p_{ki}^{(m)}>0$,
wtedy $\pi_i=(\pi P^m)_i\geq\pi_kp_{ki}^{(m)}>0$.

Dla dowolnego $i$ nie zakładamy jeszcze powracania.
Definiujemy $v_j=\sum_{n\geq0}\Prob_i(X_n=j,\tau>n)$,
dopuszczając wartość nieskończoną. W stacjonarnym
łańcuchu, dla ustalonego $M\geq0$, zdarzenia
\[
A_{n,j}^{(M)}=\{X_{M-n}=i,\ X_M=j\}
\cap\bigcap_{\ell=M-n+1}^{M}\{X_\ell\ne i\},
\qquad 0\leq n\leq M,
\]
są rozłączne względem $n$ przy ustalonym $j$.
Dla $n=0$ przecięcie jest puste w sensie
liczby czynników, czyli równe całej przestrzeni;
wtedy $A_{0,j}^{(M)}=\{X_M=i,X_M=j\}$.
Każde określa ostatnią wizytę w $i$ do chwili $M$
i zawiera się w $\{X_M=j\}$. Z prawa trajektorii
oraz stacjonarności prawdopodobieństwo tego
zdarzenia wynosi $\pi_i\Prob_i(\tau>n,X_n=j)$.
Sumując po $n$ i przechodząc do granicy, dostajemy
\[
\pi_j\geq\pi_i v_j.
\]
Sumując po $j$ i stosując twierdzenie Tonellego,
otrzymujemy
\[
1\geq\pi_i\sum_jv_j
=\pi_i\sum_{n\geq0}\Prob_i(\tau>n)
=\pi_i\E_i\tau.
\]
Zatem $\E_i\tau<\infty$, co również implikuje
$\Prob_i(\tau<\infty)=1$. Ponieważ $i$ był dowolny,
dowodzi to (3)$\Rightarrow$(2); implikacja
(2)$\Rightarrow$(1) jest oczywista.

Pozostaje ustalić jedyność i dokładny wzór.
Dla ustalonego $i$ z pierwszej części dowodu mamy
$vP=v$, $v_i=1$. Wektor $w=\pi-\pi_iv$ jest
nieujemny na mocy wykazanej nierówności,
sumowalny i spełnia $wP=w$ oraz $w_i=0$.
Dla dowolnego $k$ wybieramy $m$ z $p_{ki}^{(m)}>0$.
Wtedy $0=w_i=(wP^m)_i\geq w_kp_{ki}^{(m)}$,
więc $w_k=0$. Stąd $\pi=\pi_iv$, a sumowanie
współrzędnych daje $1=\pi_i\E_i\tau$.
Wektor $v$ zależy wyłącznie od $P$ i wybranego $i$,
więc $\pi=v/(\E_i\tau)$ wyznacza cały rozkład
stacjonarny jednoznacznie.
\end{proof}

\begin{uwaga}
W nieprzywiedlnym łańcuchu wszystkie stany są
jednego z trzech typów: przemijające, zerowo
powracające albo dodatnio powracające. Istnienie
rozkładu stacjonarnego odróżnia przypadek dodatniego
powracania od dwóch pozostałych, ale jego brak sam
w sobie nie rozstrzyga, czy łańcuch jest przemijający,
czy zerowo powracający.
\end{uwaga}

\begin{przyklad}[Spacer symetryczny jest zerowo powracający]
Wykazaliśmy wcześniej powracanie spaceru symetrycznego
na $\Z$. Gdyby istniał rozkład stacjonarny $\pi$,
dla każdego $j\in\Z$ zachodziłoby
$\pi_j=(\pi_{j-1}+\pi_{j+1})/2$.
Zatem $\pi_{j+1}-\pi_j$ nie zależy od $j$.
Ciąg $(\pi_j)_{j\in\Z}$ byłby afiniczny,
a jego nieujemność dla wszystkich całkowitych
$j$ wymusza zerową różnicę. Istotnie, gdyby różnica była dodatnia, masy
dla dostatecznie małych $j$ byłyby ujemne;
gdyby była ujemna, stałoby się tak dla
dostatecznie dużych $j$. Stały ciąg nieujemny
na $\Z$ nie może mieć sumy równej $1$.
Nie istnieje więc rozkład stacjonarny.
Z twierdzenia~\ref{thm:rp1-12-positive-recurrence}
wynika, że wszystkie stany są zerowo powracające.
\end{przyklad}

\begin{twierdzenie}[Istnienie rozkładu stacjonarnego na skończonym zbiorze]
Każda macierz stochastyczna $P$ na niepustym
skończonym zbiorze $S$ ma co najmniej jeden
rozkład stacjonarny.
\end{twierdzenie}
\begin{proof}
Wybierzmy dowolny rozkład $\mu$ i zdefiniujmy
$\nu_N=N^{-1}\sum_{n=0}^{N-1}\mu P^n$.
Są to rozkłady probabilistyczne w skończenie
wymiarowym sympleksie
$\{x_i\geq0,\ \sum_i x_i=1\}$, który jest zwarty.
Istnieje więc podciąg $N_k$ i rozkład $\pi$ taki,
że $\nu_{N_k}\to\pi$ współrzędnie. Ponieważ
\[
\nu_NP-\nu_N=\frac{\mu P^N-\mu}{N},
\]
prawa strona zbiega do zera współrzędnie
(jej współrzędne są ograniczone modułem przez $1/N$).
Przechodząc do granicy i korzystając z ciągłości
mnożenia macierzy w skończonym wymiarze,
otrzymujemy $\pi P=\pi$.
\end{proof}

\begin{twierdzenie}[Jedyność rozkładu stacjonarnego dla skończonego łańcucha]
Jeśli $S$ jest skończony, a $P$ nieprzywiedlna,
to rozkład stacjonarny $\pi$ istnieje, jest jedyny
i spełnia $\pi_i>0$ dla każdego $i\in S$.
\end{twierdzenie}
\begin{proof}
Istnienie wynika z poprzedniego twierdzenia,
a jedyność także z twierdzenia~\ref{thm:rp1-12-positive-recurrence}.
Podajmy niezależny dowód elementarny dla skończonego
$S$, który nie wymaga analizy cykli powrotu.
Jeżeli $\pi_i>0$ i $i\to j$, to dla pewnego $n$
$\pi_j=(\pi P^n)_j\geq\pi_i p_{ij}^{(n)}>0$.
Ponieważ pewne $\pi_i>0$ i łańcuch jest
nieprzywiedlny, wszystkie współrzędne są dodatnie.
Niech $\rho$ będzie innym rozkładem stacjonarnym;
z tego samego argumentu $\rho_i>0$.
Niech $c=\max_{i\in S}\pi_i/\rho_i$ i
$D_i=c\rho_i-\pi_i\geq0$. Co najmniej jeden $D_i=0$.
Ponieważ obie miary są stacjonarne,
$D_j=\sum_iD_i p_{ij}$ dla wszystkich $j$.
Wybieramy $j$ z $D_j=0$. Z $DP^n=D$ dla
każdego $n$ otrzymujemy
$0=D_j\geq D_i p_{ij}^{(n)}$. Dla każdego $i$
nieprzywiedlność zapewnia $n$ z $p_{ij}^{(n)}>0$,
stąd $D_i=0$ dla wszystkich $i$. Zatem $\pi=c\rho$.
Sumowanie współrzędnych daje $c=1$, więc
$\pi=\rho$.
\end{proof}

\begin{lemat}[Wspólna dodatnia potęga w przypadku nieokresowym]
Niech $S$ będzie skończony, $P$ nieprzywiedlna
i niech okres stanów wynosi $1$. Wtedy istnieje
$N\geq1$, dla którego wszystkie współrzędne
macierzy $P^N$ są dodatnie.
\end{lemat}
\begin{proof}
Ustalmy stan $o$. Z definicji okresu liczby
$n\geq1$ z $p_{oo}^{(n)}>0$ mają największy
wspólny dzielnik równy $1$. Istnieje zatem
\emph{skończony} podzbiór długości
$n_1,\ldots,n_r$ o tym samym NWD: kolejne NWD
liczb naturalnych tworzą ciąg nierosnący,
który po skończonej liczbie zmian stabilizuje się.
Każda dostatecznie duża liczba naturalna jest
nieujemną kombinacją całkowitą $n_1,\ldots,n_r$.
Aby to uzasadnić, oznaczmy $a=n_1$.
Reszty modulo $a$ osiągane przez nieujemne sumy
$n_1,\ldots,n_r$ tworzą podmonoid skończonej
grupy $\Z/a\Z$. Jest on podgrupą, ponieważ
dla każdego elementu $x$ tego podmonoidu istnieje
$d\geq1$ z $dx=0$, a wtedy $-x=(d-1)x$ również
należy do podmonoidu. Z tożsamości Bézouta dla
liczb $n_1,\ldots,n_r$ wynika, że podgrupa zawiera
resztę $1$, a więc jest całą grupą.
Wybierzmy dla każdej reszty $b$ jedną nieujemną sumę
$h_b$ o tej reszcie. Każda liczba o danej reszcie
większa lub równa $h_b$ jest postaci $h_b+ka$,
$k\geq0$, więc jest realizowana jako długość
ścieżki z $o$ do $o$.

Dla każdego $i,j\in S$ wybierzmy dodatnie
ścieżki $i\to o$ i $o\to j$ o długościach
$r_i,s_j$; można tu dopuścić długość zero.
Wobec poprzedniego ustępu, gdy $N-r_i-s_j$
jest dostatecznie duże, istnieje ścieżka
$i\to o\to o\to j$ długości dokładnie $N$.
Ponieważ par $(i,j)$ jest skończenie wiele,
wybierzmy wspólne duże $N$. Wtedy
$p_{ij}^{(N)}>0$ dla wszystkich par.
\end{proof}

\begin{twierdzenie}[Zbieżność skończonego łańcucha ergodycznego]
Niech $S$ będzie skończony, $P$ nieprzywiedlna
i nieokresowa. Wtedy istnieje jedyny rozkład
stacjonarny $\pi$ i dla każdego początkowego
rozkładu $\mu$ zachodzi
\[
\sum_{j\in S}|(\mu P^n)_j-\pi_j|\longrightarrow0.
\]
Co więcej, istnieją stałe $C<\infty$ i $r\in(0,1)$
(zależne od $P$), dla których lewa strona jest
nie większa niż $Cr^n$.
\end{twierdzenie}
\begin{proof}
Z poprzedniego lematu wybierzmy $N$ takie,
że $Q=P^N$ ma wszystkie współrzędne dodatnie.
Niech $m=|S|$ i $\varepsilon=m\min_{i,j}Q_{ij}>0$.
Ponieważ każdy wiersz $Q$ sumuje się do $1$,
$\varepsilon\leq1$. Gdy $\varepsilon=1$,
wszystkie wiersze $Q$ są jednostajne. Z
$\pi Q=\pi$ wynika wtedy, że $\pi$ jest
rozkładem jednostajnym i $\mu Q=\pi$ dla każdego
rozkładu $\mu$.
Rozpatrzmy teraz przypadek $0<\varepsilon<1$. Mamy
\[
Q_{ij}=\varepsilon\,u_j+(1-\varepsilon)R_{ij},
\qquad u_j=1/m,
\]
gdzie $R_{ij}=(Q_{ij}-\varepsilon/m)/(1-\varepsilon)$
jest nieujemne, a suma każdego wiersza wynosi
$(1-\varepsilon)/(1-\varepsilon)=1$. Zatem $R$
jest macierzą stochastyczną.
Dla dowolnych rozkładów $\mu,\nu$
$\mu Q-\nu Q=(1-\varepsilon)(\mu R-\nu R)$.
Dla każdej macierzy stochastycznej $K$ i dowolnego
wektora $v$ mamy
\[
\|vK\|_1=\sum_j\left|\sum_i v_iK_{ij}\right|
\leq\sum_i |v_i|\sum_jK_{ij}=\|v\|_1.
\]
Stosując to do $K=R$ oraz $v=\mu-\nu$, otrzymujemy
\[
\|\mu Q-\nu Q\|_1
\leq(1-\varepsilon)\|\mu-\nu\|_1.
\]
Ustalmy stacjonarny $\pi$ i podstawmy $\nu=\pi$.
Iteracja daje
$\|\mu Q^k-\pi\|_1\leq2(1-\varepsilon)^k$.
Dla $n=kN+s$, $0\leq s<N$, kontrakcja przez
$P^s$ daje to samo oszacowanie dla
$k=\lfloor n/N\rfloor$.
Przy $0<\varepsilon<1$ jest on geometryczny:
można wziąć $r_0=(1-\varepsilon)^{1/N}<1$
i stałą $C=2/(1-\varepsilon)$.
Przy $\varepsilon=1$ mamy $\mu P^n=\pi$
dla $n\geq N$, gdyż $\mu P^N=\pi$ i $\pi P=\pi$.
Dla $n<N$ norma różnicy rozkładów jest nie większa
niż $2$. Wystarczy więc dowolne $r_0\in(0,1)$
oraz $C=2r_0^{-(N-1)}$.
\end{proof}

\begin{uwaga}[Odległość całkowitej wariacji]
Na zbiorze przeliczalnym definiujemy
$\|\mu-\nu\|_{\mathrm{TV}}=\sup_{A\subseteq S}
|\mu(A)-\nu(A)|=\frac12\sum_i|\mu_i-\nu_i|$.
Użyta w dowodzie norma $\ell^1$ jest dwa razy
większa od odległości całkowitej wariacji.
Aby uzasadnić wzór, połóżmy $v_i=\mu_i-\nu_i$
i $A_+=\{i:v_i>0\}$. Ponieważ $\sum_i v_i=0$,
suma dodatnich współrzędnych jest równa modułowi
sumy ujemnych i wynosi $\tfrac12\sum_i|v_i|$.
Dla dowolnego $A\subseteq S$ moduł $\sum_{i\in A}v_i$
nie przekracza tej liczby; dla $A=A_+$ zachodzi
równość. Wszystkie szeregi są absolutnie zbieżne,
gdyż $\mu$ i $\nu$ są rozkładami probabilistycznymi.
W przypadku przestrzeni nieskończonych rozkład
stacjonarny może nie istnieć, a okresowość nie
jest jedyną przeszkodą zbieżności.
\end{uwaga}


\begin{przyklad}[Stacjonarność bez niezależności obserwacji]
Niech $S=\{0,1\}$, a macierz przejścia wynosi
$P=\begin{psmallmatrix}3/4&1/4\\1/4&3/4\end{psmallmatrix}$.
Rozpocznijmy proces od rozkładu stacjonarnego
$\pi=(1/2,1/2)$, co sprawdzamy bezpośrednio
z równania $\pi P=\pi$. Wtedy dla każdego $n$
$\Prob(X_n=1)=1/2$, lecz
$\Prob(X_0=1,X_1=1)=\pi_1p_{11}=3/8$,
podczas gdy iloczyn mas brzegowych wynosi $1/4$.
Zatem ciąg jest stacjonarny, ale kolejne stany
nie są niezależne. Rozkład każdej pojedynczej obserwacji
nie wyznacza zależności czasowej procesu.
\end{przyklad}

\section{Problemy wejścia jako równania harmoniczne}

\begin{twierdzenie}[Prawdopodobieństwo trafienia]
Niech $A,B\subseteq S$ będą rozłączne, a
$\tau_A,\tau_B$ czasami wejścia.
Dla $i\in S$ połóżmy
$h(i)=\Prob_i(\tau_A<\tau_B,\ \tau_A<\infty)$.
Wtedy $0\leq h\leq1$, $h=1$ na $A$,
$h=0$ na $B$ oraz
\[
h(i)=\sum_{j\in S}p_{ij}h(j)
\quad(i\notin A\cup B).
\]
Co więcej, $h$ jest \emph{najmniejszą}
nieujemną funkcją spełniającą te warunki brzegowe
i równanie na $S\setminus(A\cup B)$.
Jeśli $\Prob_i(\tau_A\wedge\tau_B<\infty)=1$
dla wszystkich $i$, to $h$ jest jedynym
ograniczonym rozwiązaniem układu z tymi
warunkami brzegowymi.
\end{twierdzenie}
\begin{proof}
Jeśli $i\in A$, mamy $\tau_A=0<\tau_B$,
a jeśli $i\in B$, to $\tau_B=0$,
więc wartości brzegowe wynoszą odpowiednio $1$ i $0$.
Dla $m\geq0$ połóżmy
$h_m(j)=\Prob_j(\tau_A\leq m,\tau_A<\tau_B)$.
Jeśli $i\notin A\cup B$, to rozdzielamy zdarzenie
trafienia w ciągu $m+1$ kroków według pierwszego kroku.
Pełna własność Markowa dla skończonego fragmentu
przyszłości daje
\[
h_{m+1}(i)=\sum_{j\in S}p_{ij}h_m(j).
\]
Nie trzeba warunkować na zdarzeniach zerowych.
Ponieważ $h_m(j)\uparrow h(j)$ dla każdego $j$,
z twierdzenia o zbieżności monotonicznej dla
szeregu nieujemnego wynika $h(i)=\sum_j p_{ij}h(j)$.

Niech $g:S\to[0,\infty)$ będzie innym rozwiązaniem
z tymi samymi wartościami brzegowymi, niekoniecznie
ograniczonym. Dla $\tau=\tau_A\wedge\tau_B$
pokażemy indukcyjnie
$g(i)=\E_i[g(X_{n\wedge\tau})]$.
Dla $n=0$ mamy $\E_i[g(X_0)]=g(i)$.
Na $\{\tau\leq n\}$ zatrzymany stan nie zmienia się
między chwilami $n$ i $n+1$. Na $\{\tau>n\}$
stan $X_n$ należy do $S\setminus(A\cup B)$.
Z własności Markowa dla zdarzeń
$\{\tau>n,X_n=k\}\in\mathcal F_n^X$ wynika
\begin{align*}
&\E_i[\chi_{\{\tau>n\}}g(X_{n+1})]\\
&\qquad=\sum_{k\notin A\cup B}\Prob_i(\tau>n,X_n=k)
                         \sum_{j\in S}p_{kj}g(j)\\
&\qquad=\sum_{k\notin A\cup B}\Prob_i(\tau>n,X_n=k)g(k)\\
&\qquad=\E_i[\chi_{\{\tau>n\}}g(X_n)].
\end{align*}
Dla nieujemnego $g$ wszystkie zamiany sumowania
uzasadnia twierdzenie Tonellego, bez odwoływania się
do warunkowej wartości oczekiwanej zmiennej, której
całkowalności jeszcze nie wykazaliśmy. Dodając
wkład zdarzenia $\{\tau\leq n\}$, dostajemy
$\E_i[g(X_{(n+1)\wedge\tau})]=\E_i[g(X_{n\wedge\tau})]$.
Indukcja dowodzi zarazem całkowalności zatrzymanych
zmiennych, gdyż $g(i)<\infty$. Ten sam argument działa
dla dowolnej rzeczywistej ograniczonej funkcji $g$,
ponieważ odpowiednie szeregi są absolutnie zbieżne.

Na $\{\tau_A\leq n,\tau_A<\tau_B\}$ mamy
$g(X_{n\wedge\tau})=1$, a poza tym z
nieujemności otrzymujemy
$g(i)\geq\Prob_i(\tau_A\leq n,\tau_A<\tau_B)$.
Zdarzenia po prawej rosną do
$\{\tau_A<\tau_B,\tau_A<\infty\}$;
ciągłość miary od dołu daje $g(i)\geq h(i)$.

Niech teraz $g:S\to\R$ będzie dowolnym ograniczonym
rozwiązaniem (niekoniecznie nieujemnym). Jeśli $\tau<\infty$
prawie na pewno przy każdym starcie, to
$X_{n\wedge\tau}=X_\tau$ od pewnego losowego
momentu, a $|g(X_{n\wedge\tau})|\leq\|g\|_\infty$.
Zbieżność zmajoryzowana prowadzi do
\[
 g(i)=\lim_n\E_i g(X_{n\wedge\tau})
 =\E_i g(X_\tau)
 =\Prob_i(\tau_A<\tau_B)=h(i).
\]
W ostatniej równości używamy wartości $g=1$
na $A$ i $g=0$ na $B$. To dowodzi jednoznaczności
w klasie funkcji ograniczonych. 
\end{proof}

\begin{przyklad}[Ruina gracza: model skończony]
Niech $S=\{0,\ldots,N\}$, gdzie $N\geq2$,
$0$ i $N$ są pochłaniające,
a dla $1\leq i\leq N-1$
przejścia $i\to i+1$ i $i\to i-1$ zachodzą
z prawdopodobieństwami odpowiednio $p\in(0,1)$
i $q=1-p$. Niech $h_i$ oznacza prawdopodobieństwo
dotarcia do $N$ przed zerem ze stanu $i$.
Równanie harmoniczne przyjmuje postać
$h_i=ph_{i+1}+qh_{i-1}$, z $h_0=0$, $h_N=1$.
Oznaczając $D_i=h_i-h_{i-1}$, otrzymujemy
$pD_{i+1}=qD_i$, więc $D_i=(q/p)^{i-1}D_1$.
Z $h_0=0$ i $h_N=1$ wynika
$h_i=D_1\sum_{k=0}^{i-1}(q/p)^k$ oraz
$D_1=(\sum_{k=0}^{N-1}(q/p)^k)^{-1}$.
Dla $p\ne q$ wzór na sumę szeregu geometrycznego daje
\[
h_i=\frac{1-(q/p)^i}{1-(q/p)^N}.
\]
Jeśli $p=q=1/2$, pierwsze różnice są stałe
i z warunku $h_N=1$ wynika $D_1=1/N$,
więc $h_i=i/N$. Wyprowadzenie nie zakłada
z góry postaci rozwiązania.

Również oczekiwany czas pochłonięcia
$u_i=\E_i(\tau_{\{0,N\}})$ można wyprowadzić
bez zgadywania. Z każdego stanu wewnętrznego
można dojść do zera przez co najwyżej $L=N-1$
kolejnych kroków w lewo. Prawdopodobieństwo tej
drogi jest co najmniej $\delta=q^{N-1}>0$.
Warunkując po każdym bloku długości $L$ na
zdarzeniu, że pochłonięcie jeszcze nie nastąpiło,
dostajemy indukcyjnie
$\Prob_i(\tau_{\{0,N\}}>kL)\leq(1-\delta)^k$.
Z tożsamości $\E_i\tau_{\{0,N\}}=
\sum_{n\geq0}\Prob_i(\tau_{\{0,N\}}>n)$ wynika
$u_i\leq L/\delta<\infty$.
Pierwszy krok i własność Markowa dają więc
$u_i=1+pu_{i+1}+qu_{i-1}$ dla $1\leq i<N$,
natomiast $u_0=u_N=0$. Przy $p=q$ rozwiązanie ma postać
$u_i=i(N-i)$, gdyż funkcja ta znika na brzegu,
a jej druga różnica wynosi
$u_{i+1}-2u_i+u_{i-1}=-2$, czyli
$u_i=1+(u_{i+1}+u_{i-1})/2$. Przy $p\ne q$, funkcja $i\mapsto i$
spełnia $p(i+1)+q(i-1)-i=p-q$,
a $h$ jest harmoniczna. Dla funkcji
$v_i=(Nh_i-i)/(p-q)$ mamy $v_0=v_N=0$ oraz
\[
pv_{i+1}+qv_{i-1}-v_i
=\frac{N(ph_{i+1}+qh_{i-1}-h_i)-(p-q)}{p-q}
=-1.
\]
Spełnia ona zatem równanie dla $u$, skąd
po zastosowaniu poniższego argumentu jedyności
\[
u_i=\frac{Nh_i-i}{p-q}.
\]
Jednoznaczność rozwiązań wynika z zasady maksimum
dla różnicy dwóch rozwiązań: dodatnie maksimum
w punkcie wewnętrznym wymusza tę samą wartość
u obu sąsiadów, ponieważ $p,q>0$, a więc
propaguje aż do brzegu, gdzie różnica jest zerowa.
Analogicznie dla ujemnego minimum rozważamy
przeciwną funkcję. Różnica musi więc być identycznie
zerowa.
\end{przyklad}

\begin{przyklad}[Łańcuch urodzeń i śmierci na skończonym odcinku]
Na $S=\{0,\ldots,N\}$ przyjmijmy
$p_{i,i+1}=b_i>0$ dla $i<N$,
$p_{i,i-1}=d_i>0$ dla $i>0$ oraz
$p_{ii}=1-b_i-d_i\geq0$, przy konwencji
$d_0=b_N=0$; wszystkie inne przejścia mają
prawdopodobieństwo zero.
Warunek równowagi szczegółowej daje
$\pi_i b_i=\pi_{i+1}d_{i+1}$.
Stąd rekurencyjnie
\[
\pi_i=C\prod_{k=0}^{i-1}\frac{b_k}{d_{k+1}},
\qquad
C^{-1}=1+\sum_{i=1}^N\prod_{k=0}^{i-1}
\frac{b_k}{d_{k+1}}.
\]
Pusty iloczyn przy $i=0$ jest równy $1$.
Wszystkie współrzędne są dodatnie, ich suma
wynosi $1$ dzięki wyborowi $C$, a równowaga
szczegółowa zachodzi dla sąsiednich stanów
i trywialnie dla pozostałych par. Łańcuch jest
nieprzywiedlny, ponieważ wszystkie przejścia
do sąsiednich stanów wewnątrz odcinka są dodatnie.
Z twierdzenia o jedyności rozkładu stacjonarnego
otrzymujemy zatem dokładnie jeden taki rozkład.
To model np. ograniczonego bufora, w którym
przyjścia zwiększają zajętość, a obsługa ją zmniejsza.
\end{przyklad}

\section{Zależność i korelacja w łańcuchu stacjonarnym}

\begin{stwierdzenie}[Kowariancja przez operator przejścia]
Niech $S$ będzie skończony, $(X_n)$ łańcuchem
o rozkładzie początkowym stacjonarnym $\pi$, a
$f,g:S\to\R$. Wtedy dla $k\geq0$
\[
\Cov(f(X_0),g(X_k))=
\sum_{i\in S}\pi_i f(i)(P^kg)(i)
-\Bigl(\sum_i\pi_i f(i)\Bigr)
 \Bigl(\sum_j\pi_jg(j)\Bigr).
\]
\end{stwierdzenie}
\begin{proof}
Warunkowanie względem $X_0$ daje
$\E[g(X_k)\mid X_0]=(P^kg)(X_0)$.
Mnożymy przez $f(X_0)$ i stosujemy wieżę.
Następnie odejmujemy iloczyn średnich,
które przy stacjonarnym rozkładzie
nie zależą od czasu.
\end{proof}

\begin{przyklad}[Autokorelacja dwustanowego łańcucha]
Niech $S=\{0,1\}$, $0<\alpha,\beta\leq1$,
$\alpha+\beta<2$ i
$P=\begin{psmallmatrix}1-\alpha&\alpha\\
\beta&1-\beta\end{psmallmatrix}$.
Przyjmijmy rozkład początkowy
$\pi=(\beta,\alpha)/(\alpha+\beta)$.
Dla $I_n=\chi_{\{X_n=1\}}$ otrzymujemy
\[
\Cov(I_0,I_k)=\pi_1(1-\pi_1)(1-\alpha-\beta)^k.
\]
Istotnie, z wcześniej wyprowadzonego wzoru
$p_{11}^{(k)}=\pi_1+(1-\pi_1)(1-\alpha-\beta)^k$.
Ponieważ $\E[I_0 I_k]=\pi_1p_{11}^{(k)}$ i
$\E I_k=\E I_0=\pi_1$, po odjęciu $\pi_1^2$
dostajemy wskazaną kowariancję. Wariancja
każdej zmiennej $I_n$ wynosi $\pi_1(1-\pi_1)>0$,
więc współczynnik autokorelacji jest równy
$(1-\alpha-\beta)^k$.
Jeśli $\alpha+\beta>1$, korelacje zmieniają
znak naprzemiennie, a jeśli $\alpha+\beta<1$,
zanikają bez zmiany znaku. Dla $\alpha+\beta=1$
korelacja jest zerowa dla każdego $k\geq1$. W tym ostatnim przypadku oba wiersze $P$
są równe $\pi$, więc prawo trajektorii jest
iloczynem mas $\pi$ i przy przyjętym starcie
zmienne są niezależne. W pozostałych przypadkach
proces ma zależność czasową mimo stałego
rozkładu jednowymiarowego.
\end{przyklad}

\section{Częstości odwiedzin i twierdzenie ergodyczne}

Istnienie rozkładu stacjonarnego nie oznacza samo
przez się, że \emph{rozkład $X_n$} musi do niego
zbiegać. Pojawia się jednak inne pytanie:
jaki odsetek czasu pojedyncza trajektoria
spędza w danym stanie? Dla skończonego
nieprzywiedlnego łańcucha odpowiedź jest
znacznie stabilniejsza i nie wymaga nieokresowości.

\begin{lemat}[Jednostajny ogon czasu wejścia w skończonym łańcuchu]
Niech $S$ będzie skończony, a $P$ nieprzywiedlna.
Dla dowolnego $i\in S$ istnieją liczby całkowite
$L\geq1$ oraz $\delta>0$ takie, że dla każdego
stanu początkowego $j\in S$ i $k\geq0$
\[
\Prob_j(\tau_i>kL)\leq(1-\delta)^k.
\]
Ponadto $\E_j\tau_i<\infty$ i
$\E_i\tau_i^+<\infty$.
\end{lemat}
\begin{proof}
Ustalmy $i\in S$. Dla każdego $j\neq i$
nieprzywiedlność zapewnia ścieżkę
$j=j_0,j_1,\ldots,j_{\ell_j}=i$ o dodatnim
prawdopodobieństwie
$a_j=\prod_{r=1}^{\ell_j}p_{j_{r-1}j_r}>0$.
Dla $j=i$ wybierzmy dodatnie przejście
z $i$ do pewnego $v$, które istnieje, gdyż
wiersz macierzy sumuje się do $1$. Jeżeli $v=i$,
mamy powrót w jednym kroku. W przeciwnym razie
dołączamy dodatnią ścieżkę z $v$ do $i$,
istniejącą dzięki nieprzywiedlności. Otrzymujemy
dodatnią ścieżkę powrotu o długości co najmniej $1$.
Ponieważ $S$ jest skończony, możemy ustalić
$L\geq\max_j\ell_j$ oraz
$\delta=\min_j a_j\in(0,1]$.
Z każdego stanu $j\neq i$ istnieje zatem
prawdopodobieństwo co najmniej $\delta$, że
$\tau_i\leq L$; ze stanu $i$ dotyczy to
$\tau_i^+$.

Niech $j\neq i$. Na zdarzeniu $\{\tau_i>kL\}$
w chwili $kL$ znajdujemy się w pewnym stanie
$X_{kL}\neq i$. Własność Markowa i oszacowanie
jednego bloku dają
\begin{align*}
\Prob_j(\tau_i>(k+1)L)
&=\E_j\Bigl[\chi_{\{\tau_i>kL\}}
  \Prob_j(\tau_i>(k+1)L\mid\mathcal F_{kL}^X)\Bigr]\\
&\leq(1-\delta)\Prob_j(\tau_i>kL).
\end{align*}
Iteracja od $k=0$ daje zadany geometryczny ogon.
Dla $j=i$ czas $\tau_i=0$, więc oszacowanie
jest trywialne.

Dla dodatniego czasu powrotu pierwszy blok
rozpoczyna się w $i$ i wykorzystujemy wybraną
ścieżkę powrotu. Jeżeli po $k\geq1$ blokach
nie było powrotu, to $X_{kL}\neq i$; od tego
momentu stosujemy drogę prowadzącą do $i$.
W obu przypadkach prawdopodobieństwo sukcesu
w następnym bloku jest co najmniej $\delta$.
Ten sam argument indukcyjny daje zatem
$\Prob_i(\tau_i^+>kL)\leq(1-\delta)^k$.
Wreszcie dla nieujemnej całkowitoliczbowej zmiennej
losowej $T$, także z dopuszczalną wartością
$\infty$, mamy punktowo
$T=\sum_{n\geq0}\chi_{\{T>n\}}$.
Twierdzenie o zbieżności monotonicznej daje
$\E T=\sum_{n\geq0}\Prob(T>n)$.
Grupując wyrazy w bloki długości $L$, otrzymujemy
\[
\E_j\tau_i\leq L\sum_{k\geq0}(1-\delta)^k
=\frac L\delta,
\qquad
\E_i\tau_i^+\leq\frac L\delta.
\]
Otrzymaliśmy jednocześnie prawie pewną skończoność
obu czasów i ich całkowalność.
\end{proof}

\begin{twierdzenie}[Częstości stanów i wzór na średni czas powrotu]
Niech $(X_n)$ będzie jednorodnym nieprzywiedlnym
łańcuchem Markowa na skończonym $S$, a $\pi$
jego jedynym rozkładem stacjonarnym. Wtedy
dla każdego $i\in S$ i dowolnego rozkładu
początkowego
\[
\frac1n\sum_{k=0}^{n-1}\chi_{\{X_k=i\}}
\longrightarrow\pi_i\quad\text{prawie na pewno}.
\]
Ponadto $\E_i\tau_i^+=1/\pi_i$, zgodnie
z twierdzeniem~\ref{thm:rp1-12-positive-recurrence}.
\end{twierdzenie}
\begin{proof}
Zacznijmy od prawa procesu $\Prob_i$, tzn.
$X_0=i$ prawie na pewno. Zdefiniujmy $T_0=0$ i
rekurencyjnie $T_r=\inf\{n>T_{r-1}:X_n=i\}$.
Z poprzedniego lematu i mocnej własności Markowa
wszystkie $T_r$ są skończone p.n. (indukcyjnie,
stosując kryterium powracania oraz wykazane wcześniej
własności kolejnych czasów powrotu).
Dla $m\geq1$ zdarzenie $L_r=T_r-T_{r-1}=m$
oznacza, że w następnych $m-1$ krokach po
$T_{r-1}$ nie odwiedzamy $i$, a w kroku $m$
odwiedzamy $i$. Jest to zdarzenie zależne od
skończonego fragmentu przyszłości. Mocna własność
Markowa daje dla każdego $m\geq1$
\[
\Prob_i(L_r=m\mid\mathcal F_{T_{r-1}}^X)
=\Prob_i(\tau_i^+=m)\quad\text{p.n.}
\]
Wcześniejsze odstępy $L_1,\ldots,L_{r-1}$ są
$\mathcal F_{T_{r-1}}^X$-mierzalne: na zdarzeniu
$\{T_{r-1}=n\}$ ich wartości są wyznaczone przez
historię do chwili $n$. Zatem dla dowolnych
$m_1,\ldots,m_r\geq1$ kolejne warunkowanie daje
\[
\Prob_i(L_1=m_1,\ldots,L_r=m_r)
=\prod_{s=1}^r\Prob_i(\tau_i^+=m_s).
\]
Dowodzi to niezależności i jednakowego rozkładu
odstępów, bez zakładania niezależności kolejnych
stanów łańcucha.
Niech $m_i=\E_i\tau_i^+\in(0,\infty)$.
MPWL (rozdział~\ref{ch:lln}) daje
\[
 \frac{T_r}{r}=\frac1r\sum_{s=1}^rL_s
 \longrightarrow m_i\quad\text{p.n.}
\]

W celu przejścia od długości cykli do liczby
odwiedzin wprowadźmy
$r(n)=\max\{r\geq0:T_r\leq n\}$.
Ponieważ wszystkie odstępy są skończone i
$T_r\geq r$, mamy $r(n)\to\infty$ przy $n\to\infty$.
Z definicji $T_{r(n)}\leq n<T_{r(n)+1}$, skąd
\[
 \frac{T_{r(n)}}{r(n)}
 \leq\frac n{r(n)}
 <\frac{T_{r(n)+1}}{r(n)+1}
   \frac{r(n)+1}{r(n)}.
\]
Dla $r(n)\geq1$ obie skrajne strony zbiegają
do $m_i$, a więc $r(n)/n\to1/m_i$.
Liczba $N_i(n)=\sum_{k=0}^{n-1}\chi_{\{X_k=i\}}$
zlicza chwile $T_r<n$, podczas gdy $r(n)$
zlicza dodatnie powroty $T_r\leq n$.
Zatem $|N_i(n)-r(n)|\leq1$, a stąd
$N_i(n)/n\to1/m_i$ prawie na pewno.

Dla dowolnego stanu początkowego $j$ mamy
$\tau_i<\infty$ p.n. (poprzedni lemat).
Z mocnej własności Markowa w $\tau_i$ wynika,
że wszystkie skończenie wymiarowe rozkłady
ciągu $(X_{\tau_i+k})_{k\geq0}$ są takie jak
dla łańcucha o prawie $\Prob_i$. Rozkład całego
ciągu jest wyznaczony przez rozkłady cylindryczne
(rozdział~\ref{ch:probability}), więc udowodniona
zbieżność częstości zachodzi również dla
przesuniętej trajektorii. Oznaczając jej liczbę
wizyt do chwili $m-1$ przez $\widetilde N_i(m)$,
dla $n>\tau_i$ mamy dokładnie
\[
\frac{N_i(n)}n=
\frac{n-\tau_i}n\,
\frac{\widetilde N_i(n-\tau_i)}{n-\tau_i}.
\]
Pierwszy czynnik dąży do $1$, a drugi do $1/m_i$
na zbiorze pełnej miary. Stąd granica
$1/m_i$ jest niezależna od stanu początkowego,
a zsumowanie po rozkładzie początkowym daje ją
dla dowolnego startu. W skończonym $S$ można
jednocześnie wybrać zbiór pełnej miary dla
wszystkich stanów $i$.

Aby rozpoznać tę granicę, rozpocznijmy łańcuch
z rozkładu stacjonarnego $\pi$. Wtedy
$\E_\pi N_i(n)/n=\pi_i$ dla każdego $n$,
bo $\Prob_\pi(X_k=i)=\pi_i$. Ponieważ
$0\leq N_i(n)/n\leq1$, zbieżność zmajoryzowana
daje $\pi_i=1/m_i$, czyli
$\E_i\tau_i^+=1/\pi_i$. Zwróćmy uwagę,
że w tym rozumowaniu \emph{nie} użyliśmy
nieokresowości: okresowy łańcuch również może
mieć stabilne częstości odwiedzin.
\end{proof}

\begin{wniosek}[Ergodyczne uśrednianie funkcji stanów]
Przy założeniach poprzedniego twierdzenia,
dla każdej funkcji $f:S\to\R$ i dowolnego
rozkładu początkowego zachodzi
\[
\frac1n\sum_{k=0}^{n-1}f(X_k)
\longrightarrow\sum_{i\in S}\pi_i f(i)
\quad\text{prawie na pewno}.
\]
\end{wniosek}
\begin{proof}
Ponieważ $S$ jest skończony,
$\frac1n\sum_{k<n}f(X_k)=
\sum_{i\in S}f(i)\frac1n\sum_{k<n}\chi_{\{X_k=i\}}$.
Dla wszystkich stanów jednocześnie zachodzi
prawie pewna zbieżność częstości, a skończona
suma pozwala przejść do granicy.
\end{proof}

\begin{przyklad}[Naprzemienne odwiedziny mimo stabilnych częstości]
Dla macierzy
$P=\begin{psmallmatrix}0&1\\1&0\end{psmallmatrix}$
i $X_0=0$ rozkład $X_n$ nie jest zbieżny:
znajduje się na przemian w stanie $0$ lub $1$.
Jednak prawie na pewno
$\frac1n\sum_{k<n}\chi_{\{X_k=0\}}\to1/2$
i analogicznie dla stanu $1$. Nieokresowość
potrzebna jest dla zbieżności \emph{rozkładów
w kolejnych chwilach}, ale nie dla \emph{częstości
na jednej długiej trajektorii}.
\end{przyklad}

\section{Macierz fundamentalna łańcucha pochłaniającego}

\begin{twierdzenie}[Liczba wizyt i prawdopodobieństwa pochłonięcia]
Niech $S=T\sqcup A$ będzie skończoną przestrzenią
stanów, gdzie $T$ i $A$ są niepuste, każdy stan
$a\in A$ jest pochłaniający (tj. $p_{aa}=1$),
a dla każdego $i\in T$ istnieje
droga o dodatnim prawdopodobieństwie do $A$.
Po uporządkowaniu stanów zapiszmy
\[
P=\begin{pmatrix}Q&R\\0&I\end{pmatrix},
\]
gdzie $Q$ opisuje przejścia wewnątrz $T$,
a $R$ z $T$ do $A$. Wtedy $Q^n\to0$, macierz
$I-Q$ jest odwracalna i
\[
N:=(I-Q)^{-1}=\sum_{n=0}^{\infty}Q^n.
\]
Współrzędna $N_{ij}$ jest oczekiwaną liczbą
wizyt w stanie $j\in T$ (łącznie z chwilą zero)
przed pochłonięciem, przy starcie w $i\in T$.
Oczekiwane czasy pochłonięcia tworzą wektor
$N\mathbf 1$, a prawdopodobieństwa pochłonięcia
w poszczególnych stanach $a\in A$ tworzą
macierz $B=NR$.
\end{twierdzenie}
\begin{proof}
Ponieważ $T$ jest skończony i z każdego stanu
istnieje dodatnia droga do $A$, znajdziemy
$L\geq1$, $\delta>0$ takie, że z dowolnego
stanu $i\in T$ prawdopodobieństwo wejścia do $A$
w najbliższych $L$ krokach jest co najmniej
$\delta$. Wybieramy dla każdego $i\in T$ jedną
dodatnią drogę do $A$, a następnie bierzemy
maksimum ich długości jako $L$ i minimum
ich prawdopodobieństw jako $\delta$.
Na zdarzeniu $\{\tau_A>kL\}$ stan $X_{kL}$
należy do $T$. Z pełnej własności Markowa
prawdopodobieństwo pozostania w $T$ przez
następne $L$ kroków, warunkowo względem
$\mathcal F_{kL}^X$, jest tam nie większe niż
$1-\delta$. Indukcyjnie otrzymujemy
$\Prob_i(\tau_A>kL)\leq(1-\delta)^k$.
Ponadto z prawa trajektorii, sumując po
stanach pośrednich w $T$, mamy
\[
(Q^n)_{ij}=\Prob_i(X_n=j,\tau_A>n),\qquad i,j\in T.
\]
Dla $n=0$ obowiązuje $Q^0=I$.
Stąd $\sum_{j\in T}(Q^n)_{ij}=\Prob_i(\tau_A>n)$.
Dla $n=kL+s$, $0\leq s<L$, ta suma nie przekracza
$(1-\delta)^k$. Zatem każdy szereg współrzędnych
macierzy $\sum_{n\geq0}Q^n$ jest zbieżny,
a $Q^n\to0$. W szczególności $\tau_A<\infty$ p.n.
Dla sumy częściowej
$(I-Q)\sum_{n=0}^{m}Q^n=I-Q^{m+1}$.
Przechodząc z $m\to\infty$, dostajemy
$(I-Q)N=I$; analogicznie $N(I-Q)=I$.

Sumując wykazaną wyżej tożsamość dla
$(Q^n)_{ij}$ po $n$ i stosując twierdzenie
o zbieżności monotonicznej, otrzymujemy
$N_{ij}=\E_i\sum_{n<\tau_A}\chi_{\{X_n=j\}}$.
Ponieważ przed $\tau_A$ w każdej chwili proces
jest w dokładnie jednym stanie z $T$, sumowanie
po $j$ daje $\E_i\tau_A=(N\mathbf 1)_i$.
Dla $a\in A$ rozbijamy zdarzenie pierwszego
wejścia do $A$ w chwili $n+1$ po stanie $j\in T$
w chwili $n$. Z własności Markowa otrzymujemy
\[
\Prob_i(\tau_A=n+1,X_{n+1}=a)
=\sum_{j\in T}(Q^n)_{ij}R_{ja}=(Q^nR)_{ia}.
\]
Zdarzenia te są rozłączne dla różnych $n$.
Sumowanie daje
$B_{ia}=\sum_{n\geq0}(Q^nR)_{ia}=(NR)_{ia}$.
Każdy wiersz $B$ sumuje się do $1$, ponieważ
pochłonięcie zachodzi prawie na pewno.
\end{proof}

\begin{przyklad}[Cztery stany, dwa cele pochłaniające]
Niech $S=\{0,1,2,3\}$, stany $0,3$ będą
pochłaniające, a przejścia z pozostałych
stanów będą następujące:
ze stanu $1$ idziemy do $0$ z prawdopodobieństwem
$1/4$ lub do $2$ z prawdopodobieństwem $3/4$;
ze stanu $2$ idziemy do $1$ z prawdopodobieństwem
$1/2$ lub do $3$ z prawdopodobieństwem $1/2$.
Dla uporządkowania $T=(1,2)$ i $A=(0,3)$ mamy
\[
Q=\begin{pmatrix}0&3/4\\1/2&0\end{pmatrix},
\qquad R=\begin{pmatrix}1/4&0\\0&1/2\end{pmatrix},
\]
a ponieważ $\det(I-Q)=1-3/8=5/8$, wzór
na macierz odwrotną daje
\[
N=\begin{pmatrix}8/5&6/5\\4/5&8/5\end{pmatrix},
\qquad B=NR=\begin{pmatrix}2/5&3/5\\1/5&4/5\end{pmatrix}.
\]
Zaczynając w stanie $1$, odwiedzamy przed pochłonięciem średnio
$8/5$ razy stan $1$ (wliczając chwilę zero)
i $6/5$ razy stan $2$,
a czas do pochłonięcia ma średnią $14/5$.
Prawdopodobieństwo zakończenia w stanie $3$
wynosi $3/5$. Zaczynając w stanie $2$, czas
ma średnią $12/5$, a szansa zakończenia w $3$
wynosi $4/5$.
\end{przyklad}


\section{Zadanie rozwiązane: trajektorie, warunkowanie i pochłanianie}

\begin{przyklad}[Urządzenie: pełne rozwiązanie jednego modelu]
Urządzenie jest w jednym z trzech stanów:
$0$ (sprawne), $1$ (awaria przejściowa),
$2$ (wyłączone na stałe). Macierz przejścia wynosi
\[
P=\begin{pmatrix}
\frac12&\frac13&\frac16\\
\frac14&\frac12&\frac14\\
0&0&1
\end{pmatrix},
\qquad
\mu=\left(\frac34,\frac14,0\right).
\]
Rozwiążemy pięć pytań, zwracając uwagę na różnicę
między warunkowaniem, sumowaniem trajektorii
i czasami zatrzymania.

\emph{(1) Jakie jest prawdopodobieństwo trajektorii
$0\to1\to0\to2$?}
Początkowy stan nie jest ustalony, dlatego
musimy uwzględnić $\mu_0=3/4$:
\[
\Prob(X_0=0,X_1=1,X_2=0,X_3=2)
=\frac34\cdot\frac13\cdot\frac14\cdot\frac16
=\frac1{96}.
\]

\emph{(2) Jak porównać przewidywanie stanu $1$
po jednej i po dwóch obserwacjach?}
Jeśli wiemy, że $X_0=0$ i $X_1=1$, własność
Markowa daje
$\Prob(X_2=1\mid X_0=0,X_1=1)=p_{11}=1/2$.
Jeżeli znamy jedynie $X_0=0$, rozważamy wszystkie
możliwe stany pośrednie:
\[
\Prob(X_2=1\mid X_0=0)
=(P^2)_{01}=\frac12\frac13+\frac13\frac12=\frac13.
\]
Te liczby są różne, ponieważ \emph{poznanie stanu
pośredniego} zmienia informację, na której opiera się prognoza.

\emph{(3) Jak odtworzyć stan początkowy z późniejszych obserwacji?}
Po zaobserwowaniu $X_1=1$ chcemy obliczyć
$\Prob(X_0=0\mid X_1=1)$. Mamy
\[
\Prob(X_1=1)=\frac34\frac13+\frac14\frac12=\frac38,
\qquad
\Prob(X_0=0\mid X_1=1)
=\frac{(3/4)(1/3)}{3/8}=\frac23.
\]
Jeśli wiemy dodatkowo, że $X_2=0$, wynik nadal
wynosi $2/3$, bo czynnik $p_{10}=1/4$ występuje
zarówno w liczniku, jak i w mianowniku. To konkretny
przykład warunkowej niezależności przeszłości
i dalszej przyszłości przy znanym stanie pośrednim.

\emph{(4) Jakie jest prawdopodobieństwo trwałego
wyłączenia w ciągu dwóch kroków, jeśli $X_0=0$?}
Stan $2$ jest pochłaniający, więc zdarzenie
$\{\tau_2\leq2\}$ jest równoważne $\{X_2=2\}$.
Można dojść do $2$ bezpośrednio w pierwszym kroku
lub w drugim kroku przez stan $0$ albo $1$:
\[
\Prob_0(\tau_2\leq2)
=\frac16+\frac12\frac16+\frac13\frac14
=\frac13.
\]
Pierwszy składnik obejmuje trajektorię, która
dotarła do stanu $2$ już w chwili $1$ i pozostała tam
do chwili $2$. Nie wolno liczyć jej ponownie jako
nowego wejścia w chwili $2$.

\emph{(5) Ile kroków średnio upłynie do trwałego wyłączenia?}
Niech $m_i=\E_i\tau_2$, gdzie $\tau_2$
jest pierwszym czasem wejścia do stanu $2$.
Z obu niepochłaniających stanów istnieje
bezpośrednie przejście do $2$ z prawdopodobieństwem
co najmniej $1/6$, więc z twierdzenia o macierzy
fundamentalnej czasy te mają skończone średnie.
Mamy $m_2=0$, a analiza pierwszego kroku daje
\[
\begin{cases}
m_0=1+\frac12m_0+\frac13m_1,\\[2pt]
m_1=1+\frac14m_0+\frac12m_1.
\end{cases}
\]
Pierwsza jedynka jest kosztem wykonanego kroku,
a po nim pozostaje oczekiwany czas dalszego
oczekiwania. Przekształcając układ otrzymujemy
$3m_0-2m_1=6$ oraz $-m_0+2m_1=4$.
Dodanie równań daje $2m_0=10$, a więc
$m_0=5$ i $m_1=9/2$.
Skończoność tych średnich jest uzasadniona:
z każdego niepochłaniającego stanu można dotrzeć
do $2$ w jednym kroku z prawdopodobieństwem
co najmniej $1/6$, więc
$\Prob_i(\tau_2>n)\leq(5/6)^n$.
Wreszcie, jeśli początkowy stan ma rozkład $\mu$,
średni czas wynosi
$\E_\mu\tau_2=(3/4)\cdot5+(1/4)\cdot(9/2)=39/8$.
\end{przyklad}

\section{Dwa dalsze modele: urna i nieskończona kolejka}

\begin{przyklad}[Model Ehrenfesta: rozkład dwumianowy]
Rozważmy $N\geq1$ rozróżnialnych cząstek,
z których każda znajduje się w jednym z dwóch
pojemników. W każdym kroku wybieramy jednostajnie
jedną cząstkę, niezależnie od dotychczasowych
wyborów, i przenosimy ją do drugiego pojemnika.
Niech $X_n\in\{0,\ldots,N\}$ oznacza liczbę
cząstek w pierwszym pojemniku. Wtedy
\[
p_{k,k+1}=\frac{N-k}{N}\quad(0\leq k<N),\qquad
p_{k,k-1}=\frac{k}{N}\quad(0<k\leq N),
\]
a wszystkie pozostałe przejścia mają
prawdopodobieństwo zero. Warunkowo względem
pełnej historii konfiguracji szanse wyboru cząstki
z obu pojemników zależą wyłącznie od $X_n$.
Własność wieży daje zatem własność Markowa
dla samego ciągu liczebności.
Warunek równowagi szczegółowej
$\pi_k(N-k)/N=\pi_{k+1}(k+1)/N$
prowadzi do rekurencji
$\pi_{k+1}=\pi_k(N-k)/(k+1)$ dla $0\leq k<N$.
Indukcyjnie $\pi_k=\pi_0\binom Nk$.
Z tożsamości $\sum_{k=0}^N\binom Nk=2^N$
otrzymujemy po normalizacji
$\pi_k=2^{-N}\binom Nk$.
Te masy spełniają równowagę szczegółową
dla wszystkich par stanów i są dodatnie.
Łańcuch jest nieprzywiedlny, ponieważ można
przejść pomiędzy dowolnymi stanami przez
kolejne przejścia do sąsiednich stanów.
Zatem jest to jedyny rozkład stacjonarny,
mianowicie rozkład dwumianowy $\Bin(N,1/2)$.

Łańcuch ma jednak okres $2$, ponieważ parzystość
liczby cząstek w pierwszym pojemniku zmienia się
w każdym kroku, a z każdego stanu istnieje
dodatnia droga powrotu długości $2$ (także
przy $N=1$). Największy wspólny dzielnik
długości powrotów jest więc równy $2$.
Rozkład $X_n$ nie zbiega z ustalonego stanu
do $\pi$: masy zbiorów stanów parzystych
i nieparzystych zmieniają się naprzemiennie
między $0$ i $1$, choć częstości odwiedzin
spełniają twierdzenie ergodyczne. Jeśli w każdym
kroku z prawdopodobieństwem $1/2$ pozostaniemy
bez ruchu, a z prawdopodobieństwem $1/2$ wykonamy
powyższe przejście, nowa macierz wynosi
$P'=(I+P)/2$. Każdy stan ma teraz dodatnią
pętlę własną, a dotychczasowe dodatnie
przejścia pozostają dodatnie, więc łańcuch
jest nieokresowy i nieprzywiedlny.
Ponieważ $\pi P'=(\pi+\pi P)/2=\pi$,
z twierdzenia o zbieżności skończonego łańcucha
ergodycznego wynika $\mu(P')^n\to\pi$ dla każdego
rozkładu początkowego $\mu$.
\end{przyklad}

\begin{przyklad}[Stabilna i niestabilna kolejka w czasie dyskretnym]
Niech $S=\N_0$, $p,q\in(0,1)$, $p+q=1$.
W stanach $i\geq1$ przechodzimy do $i+1$
z prawdopodobieństwem $p$ i do $i-1$
z prawdopodobieństwem $q$.
W stanie $0$ przechodzimy do $1$ z
prawdopodobieństwem $p$ i pozostajemy w $0$
z prawdopodobieństwem $q$. Taki łańcuch
modeluje liczbę klientów w buforze z przyjściami
i odchodzeniem pojedynczych jednostek.

Równowaga szczegółowa daje
$\pi_i p=\pi_{i+1}q$, czyli
$\pi_i=\pi_0(p/q)^i$. Jeśli $\rho=p/q<1$,
możemy unormować szereg geometryczny,
otrzymując
\[
\pi_i=(1-\rho)\rho^i,
\qquad \E_\pi X_0=\frac{\rho}{1-\rho}.
\]
Normalizacja wynika z $\sum_{i\geq0}\rho^i=(1-\rho)^{-1}$,
a wzór na średnią z
$\sum_{i\geq1}i\rho^i=\rho/(1-\rho)^2$
(różniczkujemy szereg potęgowy dla $|\rho|<1$).
Otrzymany rozkład jest stacjonarny dzięki
równowadze szczegółowej.
Aby wykazać jednoznaczność i rozstrzygnąć
przypadek $p\geq q$, załóżmy, że dowolny rozkład
$\nu$ jest stacjonarny. Dla $C_i=\{0,\ldots,i\}$
z warunku $\nu P=\nu$ otrzymujemy
\[
\nu(C_i)=\nu(C_i)-\nu_i p+\nu_{i+1}q.
\]
W istocie jedyne przejścia przez brzeg $C_i$
to $i\to i+1$ oraz $i+1\to i$.
Stąd $\nu_i p=\nu_{i+1}q$ dla każdego $i$,
czyli $\nu_i=\nu_0\rho^i$. Masa $\nu_0$ musi
być dodatnia, bo inaczej wszystkie masy
byłyby zerowe. Dla $\rho<1$ normalizacja
wyznacza ją jednoznacznie, a dla $\rho\geq1$
szereg mas jest rozbieżny. W tym drugim
przypadku nie istnieje rozkład stacjonarny.

Łańcuch jest nieprzywiedlny, gdyż $p,q>0$.
Z twierdzenia~\ref{thm:rp1-12-positive-recurrence}
wynika, że dla $p<q$ wszystkie stany są dodatnio
powracające, a w szczególności
$\E_0\tau_0^+=1/(1-\rho)$.
Aby rozróżnić przypadki $p=q$ i $p>q$,
obliczmy prawdopodobieństwo wejścia do zera
przy starcie z $1$. Dla całkowitego $M>1$
wprowadźmy $\tau_M$, czas wejścia do stanu $M$.
Do chwili $\tau_0\wedge\tau_M$ proces zachowuje
się jak spacer z przykładu ruiny gracza.
Stąd
\[
\Prob_1(\tau_0<\tau_M)=
\begin{cases}
1-1/M,&p=q,\\
1-\dfrac{1-q/p}{1-(q/p)^M},&p\ne q.
\end{cases}
\]
Zdarzenia $\{\tau_0<\tau_M\}$ rosną wraz z $M$:
aby osiągnąć $M+1$, trzeba wcześniej osiągnąć $M$.
Ich suma jest równa $\{\tau_0<\infty\}$.
Istotnie, nierówność $\tau_0<\tau_M$ implikuje
$\tau_0<\infty$. Odwrotnie, dla każdej trajektorii
z $\tau_0<\infty$ wystarczy wybrać $M$ większe
od maksimum stanów odwiedzonych do $\tau_0$,
aby zachodziło $\tau_0<\tau_M$. Z ciągłości miary z dołu
otrzymujemy
\[
\Prob_1(\tau_0<\infty)=
\begin{cases}
1,&p\leq q,\\
q/p,&p>q.
\end{cases}
\]
Rozdzielając pierwszy krok ze stanu $0$, mamy
$\Prob_0(\tau_0^+<\infty)=q+p\Prob_1(\tau_0<\infty)$:
z prawdopodobieństwem $q$ następuje powrót już
w chwili $1$, a z prawdopodobieństwem $p$
proces zaczyna dalszą drogę od stanu $1$.
Dla $p=q$ jest to $1$, lecz nie istnieje
rozkład stacjonarny, zatem łańcuch jest zerowo
powracający. Dla $p>q$ prawdopodobieństwo
powrotu wynosi $2q<1$, więc łańcuch jest
przemijający. W ten sposób otrzymujemy pełną
klasyfikację kolejki, a nie tylko warunek istnienia
rozkładu stacjonarnego.
\end{przyklad}


\section{Granice zakresu przedstawionej teorii}

Zajmowaliśmy się \emph{przeliczalnymi} przestrzeniami
stanów i \emph{dyskretnym} czasem. Dla ogólnych
przestrzeni stanów macierz zastępuje jądro
przejścia $K(x,A)$, mierzalne w $x$ i będące
miarą probabilistyczną w $A$. Przy czasie ciągłym
pojawiają się półgrupy $(P_t)_{t\geq0}$,
generatory i problemy regularności trajektorii.
Systematyczną teorię tych procesów rozwiniemy
w tomie II; tutaj konstrukcje dyskretne
przygotowują formalizm filtracji i martyngałów.


